\documentclass{article}
\usepackage[T1]{fontenc}
\usepackage[utf8]{inputenc}
\usepackage{geometry}
\geometry{margin=1.25in}
\usepackage{amsmath,amsfonts,amssymb,amsthm,thmtools,mathtools}
\usepackage{enumerate}

%\usepackage[notref,notcite,color]{showkeys}
\usepackage{graphicx}
%\usepackage{MnSymbol}
%\usepackage{lscape,array}
%\usepackage[T1]{fontenc}
%\usepackage{color}
%\renewcommand{\thefootnote}{\alph{footnote}}
%\numberwithin{equation}{section}
%\usepackage{textcomp}
%\usepackage{amsbsy}
\usepackage{comment}
\usepackage[noabbrev,capitalise]{cleveref}
\usepackage{todonotes}
%\usepackage{chngcntr}

\usepackage[backend=bibtex,style=numeric]{biblatex}
\addbibresource{hadram.bib}

\newtheorem{thm}{Theorem}[section]
\newtheorem{lem}[thm]{Lemma}
\newtheorem{obs}[thm]{Observation}
\newtheorem{prop}[thm]{Proposition}
\newtheorem{cor}[thm]{Corollary}
\newtheorem{conj}[thm]{Conjecture}
\newtheorem{fact}[thm]{Fact}
\newtheorem{claim}{Claim}
\newtheorem{remark}[thm]{Remark}

\newcommand{\la}{\langle}
\newcommand{\ra}{\rangle}
\newcommand{\sgn}{\operatorname{sgn}}
\newcommand{\eps}{\varepsilon}
\newcommand{\mf}[1]{\mathfrak{#1}} 
\newcommand{\mc}[1]{\mathcal{#1}} 
\newcommand{\bb}[1]{\mathbb{#1}}
\newcommand{\brm}[1]{\operatorname{#1}}
\newcommand{\bd}[1]{\partial #1}
%%

\newcommand{\fs}[2]{\left(\frac{#1}{#2}\right)}
\newcommand{\s}[1]{\left(#1\right)}
%%
\newcommand{\one}{\boldsymbol{1}}
\newcommand{\z}{\boldsymbol{0}}
\usepackage[noabbrev,capitalise]{cleveref}
\renewcommand{\v}{\textup{\textsf{v}}}
\newcommand{\e}{\textup{\textsf{e}}}
\renewcommand{\d}{\textup{\textsf{d}}}
\newcommand{\cX}{\mathcal X}
\newcommand{\N}{\mathbb N}
\newcommand{\Prob}{\mathbb{P}}
\DeclareMathOperator{\had}{h}
%\renewcommand{\cr}{\brm{cr}}

\definecolor{definitionred}{RGB}{128,0,32}

\newcommand{\defn}[1]{\textcolor{definitionred}{\emph{#1}}}

\crefname{thm}{Theorem}{Theorems}
\crefname{lem}{Lemma}{Lemmas}
\crefname{cor}{Corollary}{Corollaries}
\crefname{obs}{Observation}{Observations}


	
	\title{Asymptotics of the three-color Hadwiger-Ramsey numbers}
	\author{ Sergey Norin		\and Neil Rahman \thanks{Department of Mathematics and Statistics, McGill University.}
	}
	
	\begin{document}
	\maketitle
	
	\begin{abstract}
		The Hadwiger-Ramsey number $R_h(k;\ell)$ is defined as the smallest integer $n$ such that for any edge coloring of a complete graph on $n$ vertices in $\ell$ colors at least one of the color classes induces a graph with a $K_k$ minor. We  show that $$R_h(k;3) = \s{\beta^{-1} + o(1)}k \sqrt{\ln k},$$
		where $\beta^{-1} = 2/\sqrt{-\ln(10-4\sqrt6)} = 1.58149\ldots$. 
		This resolves the first open case of a problem by Axenovich and Steiner, who determined the asymptotic behavior of $R_h(k;2)$ and asked about the asymptotics of $R_h(k;\ell)$ for $\ell \geq 3$.
		\end{abstract}
	
\section{Introduction}

An \defn{$\ell$-decomposition of the complete graph $K_n$} is a collection of of $K_n$ such that $(E(G_1),\ldots,E(G_\ell))$ is a partition of $E(K_n)$. Thus $\ell$-decompositions correspond to colorings of edges of $K_n$.
The \defn{classical Ramsey number $R(k;\ell)$} is defined as the smallest integer $n$ such that for any $\ell$-decomposition $(G_1,\ldots,G_\ell)$ of $K_n$ at least one of the graphs  $G_i$ contains a $K_k$ subgraphs. Numerous generalizations of classical Ramsey numbers have been extensively investigated...

In particular, in $1974$ Chartrand and Polimeni~\cite{} have generalized classical Ramsey numbers to  $f$-Ramsey numbers, where for a monotone graph parameter $f$ the \defn{$f$-Ramsey number $R_f(k;\ell)$} is defined as the smallest integer $n$ such that for every $\ell$-decomposition $(G_1,\ldots,G_\ell)$ of $K_n$ we have $f(G_i) \geq k$ for some $1 \leq i \leq \ell$.  Thus $R(k;\ell)=R_{\omega}(k;\ell)$, where $\omega(G)$ denotes the clique number. Chartrand and Polimeni exactly determined the $\chi$-Ramsey numbers, where $\chi$ denotes the chromatic number showing that $$R_{\chi}(k;\ell)=1+(k-1)^{\ell-1}$$ for all pairs of positive integers $k$ and $\ell$.

The $f$-Ramsey numbers have been further investigated for a variety of natural parameters $f$. For example \begin{itemize} \item $R_{\chi'}(k;\ell)$ where $\chi'$ denotes the chromatic index has been completely determined by Lesniak-Foster and Roberts~\cite{},
	\item Matula~\cite{} considered Ramsey theory for graph connectivity showing $$ 2\ell(k-1)+1 \leq R_{\kappa}(k;\ell) \leq \frac{10}{3}\ell(k-1)+1,$$ where $\kappa^{\star}(G)$ denotes the largest integer $k$ such that $G$ contains an $k$-connected subgraph. \footnote{The corresponding natural parameter $\kappa$ equal to graph connectivity is not monotone.} The upper bound can be improved to $\frac{19}{6}\ell (k-1) -k +3$, using the bounds of Bernshteyn and Kostochka~\cite{} on  the maximum number of edges in a graph with no  $k$-connected subgraph. 
	
\end{itemize} 

See also...

The \defn{Hadwiger number $\had(G)$} of a graph $G$ is defined as the maximum $k$ such that $G$ has a $K_k$ minor. The $h$-Ramsey numbers, which we refer to as the \defn{Hadwiger-Ramsey numbers}, were first systematically investigated only recently by Axenovich and Steiner~\cite{hadramtwo} , who have determined the asymptotic behavior of the two-color Hadwiger-Ramsey numbers, showing that   \begin{equation}\label{e:2color}R_{h}(k;2) = (1+o(1))k\sqrt{\log_2{k}}.\end{equation} 

They have left the problem of determining the  asymptotics of $R_{h}(k;\ell)$ for fixed $\ell \geq 3$ open. Our main result resolves the first  case of this problem.

Let $$\beta = \frac{1}{2}\sqrt{-\ln(10-4\sqrt6)}$$
be the constant ubiquitous in our calculations.

\begin{thm}\label{t:main}  
	$$R_h(k;3) = \s{\beta^{-1} + o(1)}k \sqrt{\ln k}.$$
\end{thm}	

The extremal graphs for both two- and three-color Hadwiger-Ramsey numbers arise from random constructions. The result of Axenovich and Steiner implies that uniformly random colorings asymptotically minimize the largest size of a monochromatic complete minor in a $2$-edge-coloring of a complete graph. 

\begin{figure}[t]
	\centering
	\includegraphics[width=\linewidth]
	{three_colour_partition_construction.pdf}
	\caption{The optimal construction for $3$-color Hadwiger-Ramsey numbers.}
	\label{f:construction}
\end{figure}

One might expect the analogous situation for $3$-edge-coloring, but it turns out that there is a better random construction. We construct $3$-decompositions of $(G_1,G_2,G_3)$ of $K_n$ such that $\had(G_i) < k$ for $n=\s{\beta^{-1} + o(1)}k \sqrt{\log k}$, certifying the lower bound on $R_h(k;3)$ in \cref{t:main}, as follows. 

Let $(X,Y)$ be a partition of the vertex set of $K_n$ into two almost equal parts and let $p=5-2\sqrt{6}$. Every edge with both ends in $X$ or both ends in $Y$ is assigned to $G_1$ with probability $1-2p$ and to each of $G_2$ and $G_3$ with probability $p$, independently. Every edge with one end in $X$ and the other in $Y$ is assigned to $G_2$ with probability $1/2$ and, otherwise, to $G_3$. See \cref{f:construction}.

 It immediately follows from our results, although we don't make it explicit, that every $3$-decomposition that is close to optimal for $3$-color Hadwiger-Ramsey numbers is close to the one described above. In particular, the proportion of edges in one of the graph in the decomposition is close to $\frac{1-2p}{2} \sim 0.399 \neq 1/3$. Thus the parts in optimal decompositions are not symmetric. This is not  unique among known Ramsey constructions, but somewhat unusual. 


\subsection{Tools and the proof outline}

The upper and lower bounds in \cref{t:main} correspond to the following two theorems.

\begin{thm}\label{existencethm} For every $\eps > 0$ there exists $n_0$ such that for every  integer $n \geq n_0$  there exists a  $3$-decomposition $(G_1, G_2,G_3)$  of $K_n$ such that 
\begin{equation}\label{e:lower} \had(G_i) \leq (1+ \eps) \beta \frac{n}{\sqrt{\log n}}, \end{equation}  for every $i \in \{1,2,3\}$.
\end{thm}

\begin{thm}\label{mainthm} For every $\eps > 0$ there exists $n_0$ such that for every  integer $n \geq n_0$ and  every $3$-decomposition $(G_1, G_2,G_3)$  of $K_n$ there exists $i \in \{1,2,3\}$  such that $$ \had(G_i) \geq (1-\eps)\beta\frac{n}{\sqrt{\log n}}. $$
\end{thm}

It remains to prove Theorems~\ref{mainthm} and~\ref{existencethm}. 

 Theorem~\ref{existencethm}  is easier. A suitable random $3$-decomposition was described in the introduction and we just need to verify that \eqref{e:lower} holds for every part of it with high probability. The techniques for estimating Hadwiger number of random graphs were first developed by Bollobas, Caitlin and Erd\H{o}s~\cite{}. While their result doesn't directly apply in our setting, the techniques extend to the following technical theorem which applies.

\begin{thm}\label{thomasonlower} For every $\eps > 0$ there exists $n_0=n_{0,\ref{thomasonlower}}(\eps)$ such that for all $n \geq n_0$ the following holds. Let $(X_1, \ldots, X_k)$ be a partition of $[n]$, let $0 \leq p_{ij}=p_{ji} \leq 1 -\eps$ for all $1 \leq i,j \leq k$. Let $G$ be a random graph with $V(G)=[n]$ defined by adding each edge $uv$ with $u \in X_i$ and $v \in X_j$ independently with probability $p_{ij}$. Let $\alpha_i = |X_i|/n$ and let
 	$$L = \sqrt{-\sum_{(i,j) \in [k]^2} \alpha_i\alpha_j\log\s{1-p_{ij}}}.$$
 	If $L > \eps$, then 
 \begin{equation}\label{e:tlower}\Prob \s{\had(G) \leq (1 +\eps) L\frac{n}{\sqrt{\log n}}} \geq 1 -\eps.\end{equation}
\end{thm}

The proof of the upper bound, \cref{mainthm}, extends the argument of Axenovich and Steiner~\cite{2603.10510} on $2$-decompositions $(G_1,G_2)$ of $K_n$. They show that if  $\had(G_1),\had(G_2) \leq O\fs{n}{\sqrt{\log n}}$  then almost all edges of $G_i$ belong to a $(\delta n)$-connected subgraph of $G_i$ for sufficiently small $\delta$.  They then apply a result of Thomason that essentially states that among all $(\delta n)$-connected graphs on $n$ vertices with given density the random graph minimizes the Hadwiger number.

Neither step directly goes through in our application.  We start with a weakening of the first step of Axenovich-Steiner~\cite{2603.10510}\todo{review citation} argument and show that almost all edges of $G_i$ belong to a union of a collection of vertex disjoint  $(\delta n)$-connected subgraphs of $G_i$. We then deduce that if none of the parts has exceedingly large Hadwiger number then for two out of three parts we can after all assume that a single $(\delta n)$-connected subgraphs contains almost all of the edges. To lower bound $\had(G_i)$, we need the following extension of the result of Thomason to a setting of \cref{thomasonlower}, where density is not fixed globally, but with respect to a vertex partition. 

\begin{thm}\label{thomasonupper} For all $\eps > 0$ there exists $n_0=n_{0,\ref{thomasonupper}}(\eps)$ such that for all $n \geq n_0$, all graphs $G$ with $\v(G) = n$ and $\kappa(G) \geq \eps n$, and all partitions $(X_1,\ldots,X_k)$ of $V(G)$ the following holds. Let $p_{ij}$  denote the density of edges between $X_i$ and $X_j$ (in particular, $p_{ii}  = \e(G[X_i])/\binom{|X_i|}{2}$) and let $\alpha_i = |X_i|/n$. Let $p_{ij}^* = \min \{p_{ij}, 1 - \eps\}$. Then $$\had(G) \geq \frac{n}{\sqrt{\log n}}\s{\sqrt{-\sum_{(i,j) \in [k]^2} \alpha_i\alpha_j\log\s{1-p^*_{ij}}} -\eps}  $$
\end{thm}

The extension is once again technical, but routine. \cref{thomasonupper} reduces the proof of \cref{mainthm} to an optimization problem. This problem, while unappealing at a first glance, allows a short solution.\todo{the $- \eps$ seems to be in the wrong place for 1.5}

\subsection*{AI disclosure}
As mentioned above, the proofs of Theorems~\ref{thomasonlower} and~\ref{thomasonupper} are technical extensions of existing proofs in the literature. They were produced by AI following an outline given by the authors, who subsequently verified and polished them. To separate the AI assisted part of the argument from the rest, we relegate these proofs to the appendix.

The remainder of the argument was initially produced and written down without AI assistance. At the final stage of preparation of the paper, AI was used for proofreading, literature search and creating the figures. 

\section{Proof of \cref{existencethm}}

In this section, we derive \cref{existencethm} from \cref{thomasonlower}. 

Clearly, it suffices to prove the theorem for even $n$, and so we assume that $n$ is even for the remainder of the proof.
Let $p=5-2\sqrt{6}$.
We assume without loss of generality that $\eps \leq p$ and let $n_0$ be chosen so that $n_0 \geq  2n_{0,\ref{thomasonlower}}(\frac{\eps}{2})$ where $n_{0,\ref{thomasonlower}}(\frac{\eps}{2})$ is as in \cref{thomasonlower}, and so that  \begin{equation}\label{e:logn2}\s{1+\frac{\eps}{2}}\frac{\sqrt{\log n}}{\sqrt{\log (n/2)}} \leq 1 + \eps\end{equation}
for all $n \geq n_0$. 

Construct a random $3$-decomposition $(G_1,G_2,G_3)$ of $K_n$ as described in the introduction. Let $(X,Y)$ be a partition of the vertex set of $K_n$ with $|X|,|Y| = \frac{n}{2}$. Every edge with both ends in $X$ or both ends in $Y$ is assigned to $G_1$ with probability $1-2p$ and to each of $G_2$ and $G_3$ with probability $p$, independently. Every edge with one end in $X$ and the other in $Y$ is assigned to $G_2$ with probability $1/2$ and, otherwise, to $G_3$.  It suffices to show \begin{equation}\label{e:t2}
	\Prob \s{\had(G_i) > (1+\eps)\beta\frac{n}{\sqrt{\log n} }} < \frac{1}{3}	
	\end{equation}
for $i=1,2,3$. 

By the choice of $n_0$ and of the decomposition, the random graph $G_1[X]$ satisfies the conclusion \eqref{e:tlower} with $k=1,\alpha_1 = 1,$ $n/2$ replacing $n$, $\eps/2$ replacing $\eps$ and $p_{11}=1-2p$, i.e.$$\Prob\s{\had(G_1[X]) > \s{1+\frac{\eps}{2}}\sqrt{-\log(2p)} \frac{n}{2\sqrt{\log (n/2)}} } \leq \frac{\eps}{2}.$$
As $\sqrt{-\log(2p)}=2\beta$ and \eqref{e:logn2} holds this implies
$$\Prob\s{\had(G_1[X]) > (1+\eps)\beta\frac{n}{\sqrt{\log n} }} < \frac{\eps}{2}.$$
By symmetry the same bound holds for $\had(G_1[Y])$. As $G_1$ is a disjoint union of $G_1[X]$ and $G_1[Y]$, we have $\had(G_1) = \max\s{\had(G_1[X]),\had(G_1[Y])}$, and thus
by the union bound \eqref{e:t2} holds for $i=1$.

It remains to verify \eqref{e:t2} for $i=2,3$. The random graph $G_2$ satisfies  the conclusion \eqref{e:tlower} with $k=2,\alpha_1 = \alpha_2=1/2, p_{11}=p_{22}=p, p_{12}=1/2$, and $\eps/2$ replacing $\eps$. Thus\todo{I think we can just use $\eps$? also n / rootlogn might be flipped} 
$$\Prob\s{\had(G_2) > \s{1+\frac{\eps}{2}}\sqrt{-\frac{1}{2}\log(1-p) -\frac{1}{2}\log \frac{1}{2}}\frac{\sqrt{\log n}}{n}} < \frac{\eps}{2}.$$
As $\sqrt{-\frac{1}{2}\log(1-p) -\frac{1}{2}\log \frac{1}{2}}=\beta$, and $\eps/2 \leq 1/3$, this implies that \eqref{e:t2} holds for $i=2$. The case $i=3$ is symmetric.




\section{Proof of \cref{mainthm}.}

Our first major goal is \cref{connectedsets} which shows that for  every $n$-vertex graph we can delete a small fraction  of  vertices,  and  of  $n^2$ edges so that the remaining graph is a disjoint union of pieces with connectivity linear in  $n$. (The actual statement is a bit more restrictive). This lemma is then used in several ways...

\begin{lem}\label[lem]{bigset}
Let $\mathcal{P}_1, \mathcal{P}_2$ be partitions of the same $n$-element set  into $\leq k$ parts. Then for every $\eps >0$ either

\begin{itemize}\item For some $i \in \{1, 2\}$, there exists $P' \in \mathcal{P}_i$ such that $|P'| \geq (1 - \eps)n$, or
\item there exist $A, B$ such that $|A|, |B| \geq \frac\eps{k^2} n$, and for all $i \in \{1,2\}$, there exist $P_{A,i}, P_{B,i} \in \mathcal{P}_i$ such that $A \subseteq P_{A,i}$, $B \subseteq P_{B,i}$, and $P_{A,i} \neq P_{B,i}$ \end{itemize}
\end{lem}%\todo{if we want to be tight, we can replace $k^2$ in the denominator with $k(k-1)$}

\begin{proof} Suppose the second condition does not hold. If we had two pairs $(P_1, Q_1), (P_2, Q_2) \in \mathcal{P}_1 \times \mathcal{P}_2$ with $|P_i \cap Q_i| \geq \frac\eps{k^2} n$ for $i \in \{1, 2\}$, $P_1 \neq P_2$, $Q_1 \neq Q_2$, we could choose $A = P_1 \cap Q_1$ and $B = P_2 \cap Q_2$ to fit the second condition. Therefore, without loss of generality there exists $P' \in \mc{P}_1$ such that for all $(P, Q) \in \mathcal{P}_1 \times \mathcal{P}_2$ with $|P \cap Q| \geq \frac\eps{k^2} n$, we have $P = P'$. But then, for all $P \in \mathcal{P}_1 \setminus \{P'\}$, we have $$|P| \leq \sum_{Q \in \mathcal{P}_2} |P \cap Q| \leq k  \cdot \frac\eps{k^2}n = \frac\eps k n.$$ Since $|\mathcal{P}_1| \leq k$  it follows that $$n - |P'|  = \sum_{P \in \mathcal{P}_1 \setminus \{P'\}} |P|\leq (k-1) \cdot \frac\eps k n \leq \eps n,$$ and therefore  $|P'| \geq (1-\eps)n$.
\end{proof}

Let $G$ be a graph. For $\eps,\eta,\delta >0$ and an integer $k \geq 0$
we say that a partition $(V_1, \ldots, V_k, X, Y)$ of $V(G)$  is an \emph{$(\eps,\eta,\delta)$-controlled partition of $G$} if \begin{itemize}
	\item[(a)] $|X| \leq \eps n$,
	\item[(b)] $G[V_i]$ is $(\delta n)$-connected for every $1 \leq i \leq k$,
	\item[(c)] $G[Y]$ is $(3\eps n)$-degenerate,
	\item[(d)] For all $v \in Y$, $v$ has are most $\eta n$ neighbours in $\bigcup_{i=1}^k V_i$,
	\item[(e)] $|V_i| \geq \eps n$ for every $1 \leq i \leq k$,
	\item[(f)] $E(V_i, V_j) = \emptyset$ for all $i \neq j$.
\end{itemize}
Note that if $(V_1, \ldots, V_k, X, Y)$ is a partition of $V(G)$ satisfying (e) then $k \leq \eps^{-1}$. In particular this bound on $k$ holds for every $(\eps,\eta,\delta)$-controlled partition.


\begin{lem}\label[lem]{connectedsets} For all $\eps, \eta > 0$, there exists $\delta = \delta_{\ref{connectedsets}}(\eps,\eta) >0$ such that every graph $G$ admits an $(\eps,\eta,\delta)$-controlled partition.
\end{lem}

\begin{proof}  We assume that $G$ is not $(3\eps n)$-degenerate, as otherwise, the partition $(\emptyset, V(G))$ is $(\eps,\eta,\delta)$-controlled for any choices of $\eta$ and $\delta$. In particular, we assume  $\eps \leq 1/3$.  Let $k  =\lfloor \eps^{-1} \rfloor$. We show that $$\delta = \min\s{\frac{\eps^3}{k!},\frac{\eps\eta}{k!}}$$ satisfies the lemma. 
	
We start by  considering a preliminary partition $(U_1,\ldots, U_{k}, X', Y)$ of $V(G)$ such that \begin{itemize}
	\item [(a')] $|X'| \leq k\eps^{3}n$,
	\item [(b')] $G[U_i]$ is $(\eps^3 n)$-connected for every $1 \leq i \leq k$,
	\item [(c)] $G[Y]$ is $(3\eps n)$-degenerate,
	\item [(e')] $U_i \neq \emptyset$ and $\deg_{G[U_i]}(v) \geq 3\eps n - k\eps^3 n$ for every $1 \leq i \leq k$ and $v \in V(G[U_i])$, and 
	\item[(k)] $k\leq \eps^{-1}$.
 \end{itemize}		
Let us show that it exists. Let $H$ be a maximal subgraph of $G$ with minimum degree $3 \eps n$, then let $Y = V(G) \setminus V(H)$. $G[Y]$ is $(3 \eps n)$-degenerate, as otherwise, there would be some $Y' \subseteq Y$ such that $H \cup G[Y']$ contradicts the choice of $H$.  Therefore the partition $(V(G)-Y, \emptyset, Y)$ satisfies (a'), (c) and (e') with $k=1$.

Choose a partition    $(U_1,\ldots, U_{k'}, X', Y)$ of $V(G)$ satisfying (a'), (c), (e') and (k) with $k'$ maximum. 
We claim that this partition satisfies (b'). 

Suppose not. Then without loss of generality there exists a proper separation $(W,W')$ of $U_{k'}$ with $|W \cap W'| \leq \eps^3 n$. But in such a case the partition  $(U_1,\ldots, U_{k'-1}, W' -W, W-W',  X' \cup (W \cap W'), Y)$ of $V(G)$ contradicts our choice of partition. Indeed, it is easy to see that conditions (a')--(e') are satisfied for the new partition. To verify (k) for the new partition, we need to show that $k' \leq \eps^{-1}-1$. But, the conditions (e') and (k) for the current partition imply that $$|U_i| \geq  3\eps n - k \eps^3 n \geq 3\eps n - \eps^2 n  \geq 2 \eps n$$ for every $i$, and so $|V(G)| \geq \sum_{i=1}^k|U_i| \geq 2\eps k n$, implying $k \leq \frac{1}{2}\eps^{-1} \leq \eps^{-1}-1,$ where the l as needed. \textbf{\textcolor{blue}{(l?)}}

This finishes the construction of the preliminary partition. 
We now replace the conditions (b') and (e')  by  
 \begin{itemize}
 	\item [(b'')] $G[U_i]$ is $(k! \delta n)$-connected for every $1 \leq i \leq k$,
 	\item [(e'')] $|U_i| \geq 2 \eps n$ for   every $1 \leq i \leq k$.
 \end{itemize}	Note that (e'') holds for the current preliminary partition, as shown above, and  $k! \delta n \leq \eps^3 n$ by (k) and the choice of $\delta$ so (b') implies (b''). 

We now  choose  a partition $(U_1,\ldots, U_{k}, X', Y)$ of $V(G)$ satisfying  (a'),(b''),(c) and (e'') with $|Y|$ minimum. (Note that (k) is implied by (e'') and so also still holds.) Then  \begin{itemize}
    \item[(d')] For all $v \in Y$, $v$ has are most $\eta n$ neighbours in $\bigcup_{i=1}^k U_i$,
\end{itemize} holds. Indeed, if  some $v \in Y$, $v$ has at least  $\eta n$ neighbours in $\bigcup_{i=1}^k U_i$, then it has at least $\eps \eta  n$ neighbours in some $U_i$. As $\eps \eta \geq k! \delta$, replacing $Y$ by $Y - \{v\}$ and $U_i$ by $U_i \cup \{v\}$ in this case contradicts minimality of $|Y|$.

Thus we can  choose  a partition $(U_1,\ldots, U_{k}, X', Y)$ of $V(G)$ satisfying  (a'),(b''),(c),(d') and (e'') and we choose it with $k$ minimum. 
Then $G[U_i\cup U_j]$ is not  $((k-1)! \delta n)$-connected  for all $1 \leq i, j \leq k$, $i \neq j$  as otherwise we would replace $U_i$ and $U_j$ by $U_i \cup U_j$ in the partition. It follows that for all such pairs $i,j$ there exists $Z_{ij} \subseteq U_i$ with $|Z_{ij}| \leq (k-1)!\delta n$ such that $E(U_i-Z_{ij},U_j)=\emptyset$.  

Finally, define\todo{isn't the convention for $\cup$ and $\bigcup$ over iterated unions the other way around}
$$V_i = U_i - \cup_{j \neq i}Z_{ij} \qquad \mathrm{and} \qquad  X = X' \bigcup \s{\cup_{\substack{1 \leq i,j \leq k \\ i \neq j}} Z_{ij}}.$$
We have $\sum_{j \neq i}|Z_{ij}| \leq (k-1)(k-1)! \delta n$. Thus $$|V_i| \geq |U_i|-(k-1)(k-1)! \delta n \geq 2\eps n-(k-1)(k-1)! \delta n \geq \eps n.$$  As $G[U_i]$ was $(k!\delta n)$-connected,  the graph $G[V_i]$ is $((k! - (k-1)(k-1)!)\delta n)$-connected and so is $(\delta n)$-connected. Thus (b) and (e) hold. As $$|X| \leq |X'| + k^2(k-1)!\delta n \leq k \eps^3n + k^2(k-1)!\delta n \leq \eps n,$$ (a) holds. The condition (c) already holds, (d') implies (d), and (f) holds by the choice of the sets $Z_{ij}$. We conclude that  $(V_1, \ldots, V_k, X, Y)$ is $(\eps,\eta,\delta)$-controlled . 
\end{proof}

Before proceeding to an application of  \cref{connectedsets}, we present two technical consequences of \cref{thomasonupper}.

\begin{lem}\label[lem]{bipMinor} For every $K > 0$ there exists $\gamma=\gamma_{\ref{bipMinor}}(K), n_0=n_{0,\ref{bipMinor}}(K) >0$ satisfying the following.  
	Let $H$ be a bipartite graph with bipartition $(A,B)$ with $|A|=|B|> n_0$ and $|E(H)|\geq (1-\gamma)|A|^2$. Then $$\had(H) \geq K \frac{|A|}{\sqrt{\log{|A|}}}.$$
	
\end{lem}
\begin{proof} We assume $K \geq 1$ without loss of generality. 
	Let $\eps$ be chosen so that $\sqrt{-\log(\eps)} \geq 12 K.$ We show that $\gamma = \frac{\eps}{2} \leq 0.01$ and let $n_0 = \max \s{ n_{0,\ref{thomasonupper}}(\eps),100}$ where $n_{0,\ref{thomasonupper}}(\eps)$ is as in \cref{thomasonupper}.
		 
Let $a = |A|$. Repeatedly delete vertices of $H$ with degree at most $\frac34a$. Suppose we perform $2c$ such deletions for some $c \leq \frac14 a$. We obtain a subgraph of $H$ with $2(a - c)$ vertices and at least $(1-\gamma)a^2-\frac32 ac$ edges. As every $2(a-c)$ vertex subgraph of $H$ has at most $(a-c)^2$ edges we have $(1-\gamma)a^2 -\frac{3}{2}ac \leq (a-c)^2$, i. e.
$$\frac12 ac - \gamma a^2 \leq c^2$$
$$\Rightarrow (c-\frac14a)^2 \geq (\frac1{16} - \gamma)a^2.$$
Since $c \leq \frac14a$ and $\gamma \leq 0.01$,
$$\frac14a - c \geq \sqrt{\frac1{16}-\gamma} a$$
Given $\gamma \leq 0.01$, we must have $c \leq \frac{1}{10}{a}$. It follows that the process must stop before arriving at the null graph, indeed otherwise the process jumps from at most $\frac{1}{5}a$ deleted vertices to more than $\frac14a$ deleted vertices in two deletions, implying $\frac1{20} a \leq 2$, contradicting $a \geq n_0 \geq 100$.
	
The graph $H'$ obtained after repeated deletions has bipartition $(A',B')$ such that $ |A'|,|B'| \leq a$,  minimum degree at least $\frac{3}{4}a$ and at least $|A'||B'|-\gamma a^2$ edges. In particular, $\kappa(H') \geq \frac{3}{4}a \geq \eps |V(H')|$ and $$\frac{|E(H')|}{|A'||B'|} \geq  1 - \frac{\gamma a^2}{\s{\frac{3}{4}a}^2} \geq 1 -2 \gamma.$$ We apply \cref{thomasonupper} to the partition $(A',B')$ of the vertex set of $H'$ with $$p^*_{12}=\min\s{\frac{|E(H')|}{|A'||B'|}, 1-\eps} =1-\eps$$ to obtain
\begin{align*}\had(H) &\geq \had(H') \geq \frac{|V(H')|}{\sqrt{\log |V(H')|}}\left(\sqrt{-\frac{|A'||B'|}{(|A'|+|B'|)^2} \log(\eps)}-\eps\right)\\  &\geq \frac{1}{2}\frac{|A|}{2\sqrt{\log {|A|}}}\sqrt{-\frac{1}{9}\log(\eps)} \geq \frac{1}{12}\frac{|A|}{\sqrt{\log {|A|}}}\cdot 12K = K \frac{|A|}{\sqrt{\log{|A|}}}.  \end{align*} 
\end{proof}	

\begin{lem}\label[lem]{denseMinor} For every $K > 0$ there exists $\gamma=\gamma_{\ref{denseMinor}}(K), n_0=n_{0,\ref{denseMinor}}(K) >0$ satisfying the following.  
Let $H$ be a graph with $|V(H)| \geq n_0$ and $|E(H)|\geq (1-\gamma)\binom{|V(H)|}{2}$. Then $$\had(H) \geq K \frac{|V(H)|}{\sqrt{\log{|V(H)|}}}.$$
\end{lem}

\begin{proof} Let $\gamma= \gamma_{\ref{bipMinor}}(3K)/2$ and let $n_0 \geq 3n_{0,\ref{bipMinor}}(3K)$ be chosen so that $2(1-\gamma)\frac{n_0-1}{n_0} \geq 1 - 2\gamma$.
	Let $a = \lfloor\frac{n_0}{2}\rfloor$. Let a pair of disjoint $A,B \subseteq V(H)$  with $|A|=|B|=a$ be chosen uniformly at random. Then $\bb{E}|E_H(A,B)| \geq \frac{2a^2|E(H)|}{|V(H)|^2} \geq (1-2\gamma)a^2.$ Fix a choice $(A,B)$ with $|E_H(A,B)| \geq (1-2\gamma)a^2$. By the choice of $\gamma$ and $n_0$ the conclusion of \cref{bipMinor} with $3K$ in place of $K$ holds for the bipartite subgraph $H'$ of $H$ with bipartition $(A,B)$ and $E(H')=E_H(A,B)$, and so we have $$\had(H) \geq \had(H') \geq 3K \frac{a}{\sqrt{\log{a}}} \geq K \frac{|V(H)|}{\sqrt{\log{|V(H)|}}},$$
	as desired.
\end{proof}	


The following is our first application of \cref{connectedsets} which also uses \cref{bigset}. It shows that if a $3$-decomposition $(G_1,G_2,G_3)$ of a complete graph $K_n$ is not too far from being extremal for Hadwiger-Ramsey numbers then two out of three graphs in the decompositions are either very sparse or have connectivity linear in $n$ after possibly deleting small exceptional set of vertices.                    

\begin{lem}\label[lem]{twobigsets} For all $\eps ,C > 0$ there exist $n_0 = n_{0,\ref{twobigsets}}(\eps,C), \delta = \delta_{\ref{twobigsets}}(\eps,C) > 0$ such that the following holds.
	
Let $(G_1,G_2,G_3)$ be a  $3$-decomposition $(G_1,G_2,G_3)$ of a complete graph $K_n$ for $n \geq n_0$.  Then either 
\begin{itemize}
\item There exist two distinct $i_1,i_2 \in [3]$ such that for each $j \in [2]$ \begin{itemize} \item either  $|E(G_{i_j})| \leq 3\eps n^2$, \item or there exists $P_j \subseteq V(G)$ such that $|P_j| \geq (1-\eps)n$ and $G_{i_j}[P_j]$ is $(\delta n)$-connected,\end{itemize} or
%\todo{by changing $\eps$, don't need complicated $1-\eps^4$, revisit}
\item  $\had(G_i) \geq C\frac n{\sqrt{\log n}}$ for some $i \in [3]$.
\end{itemize}
\end{lem}

\begin{proof} Let $0<\eps'<\frac1{10}$ be chosen so that 
$
(1-\eps')(1-2\eps') \geq (1-\eps)$ and $\eps' \leq \frac{2}{3}\eps$.
 Let $K=2C(\eps')^{-3}$. Let $\gamma=\gamma_{\ref{bipMinor}}(K)$ and $n'_0=n_{0,\ref{bipMinor}}(K)$ be as in \cref{bipMinor}. Let $\eta =\frac{1}{4}(\eps')^{3}\gamma$.
By \cref{connectedsets} there exists $\delta > 0$ such that every graph admits an $(\eps',\eta,\delta)$-controlled partition. Finally, let $n_0 = 2 (\eps')^3n'_0$. We show that these $\delta$ and $n_0$ satisfy the lemma.

Assume that the first outcome of the lemma does not hold, i.e. we assume without loss of generality that $|E(G_{i})| \geq 3\eps n^2$  and there does not exist $P_i \subseteq V(G)$ such that $|P_i| \geq (1-\eps)n$ and $G_{i}[P_i]$ is $(\delta n)$-connected for $i=1,2$. 
Let 
$(V_{1,i}, \ldots, V_{k_i,i}, X_i, Y_i)$ be  an $(\eps',\eta,\delta)$-controlled partition of $G_i$. As $G_i[V_{j,i}]$ is $(\delta n)$-connected, we have $|V_{j,i}| \leq (1-\eps)n$ for  $i=1,2$ and $j \in [k_i]$. Moreover, $$3\eps n^2 \leq |E(G_i)| \leq |E(G_i[Y_i])| + |V(G)-Y_i| n \leq 3\eps'n^2 + |V(G)-Y_i| n,$$ implying $|V(G)-Y_i| \geq 3(\eps -\eps')n \geq \eps n$

Apply \cref{bigset} with $\eps'$ in place of $\eps$ and $k =\max(k_1,k_2)+1$ to the pair of  partitions
$$\mathcal{P}_1 = (V_{1,1} \setminus X_2, \ldots, V_{k_1,1} \setminus X_2, Y_1 \setminus X_2),$$
$$\mathcal{P}_2 = (V_{1,2} \setminus X_1, \ldots, V_{k_2,2} \setminus X_1, Y_2 \setminus X_1)$$
of $V(G) \setminus X_1 \setminus X_2$. 
As we have shown above, each part of $\mc{P}_i$ has size at most $$(1-\eps)n \leq (1-\eps')(1-2\eps')n \leq (1-\eps')|V(G) \setminus X_1 \setminus X_2|.$$ Thus the first outcome of \cref{bigset} cannot hold.
Further note that $k \leq (\eps')^{-1}+1$.
 So there exist $A, B \subseteq V(G)$ such that
$$|A|, |B| \geq \frac{\eps'}{(k+1)^2}(1-2\eps') n \geq \frac{\eps'(1-2\eps')}{((\eps')^{-1}+2)^2} n\geq \frac{1}{2}(\eps')^3n$$
and for all $i \in \{1,2\}$, there exist $P_{A,i}, P_{B,i} \in \mathcal{P}_i$ such that $A \subseteq P_{A,i}$, $B \subseteq P_{B,i}$, and $P_{A,i} \neq P_{B,i}$. Without loss of generality we assume $|A|=|B|$. Let $a = |A|$. 
By definition of the $(\eps',\eta,\delta)$-controlled partition it follows that  either every vertex in $P_{A,i}$ has at most $\eta n$ neighbours in $P_{B,i}$ or every vertex of 
 $P_{B,i}$ has at most $\eta n$ neighbors in $P_{A,i}$. Thus there are at most $\eta n a$ edges between $A$ and $B$ in each of $G_1$ and $G_2$. 
 
Consider the bipartite subgraph $H$ of $G_3$ with $V(H)=A \cup B$ and $E(H)=E_{G_3}(A,B)$. We have $$|E(H)| \geq a^2 - 2\eta n a \geq (1- 4(\eps')^{-3}\eta) a^2 = (1-\gamma)a^2$$ and $a \geq \frac{1}{2}(\eps')^3 n \geq n'_0$. By \cref{bipMinor} $$\had(H) \geq K\frac{ a}{\sqrt{\log a}} \geq  K\frac {\frac{1}{2}(\eps')^3n}{\sqrt{\log n}} \geq C\frac{ n}{\sqrt{\log n}},$$
and so the second outcome of the lemma holds with $i=3$.
\end{proof}

We need a couple of additional optimization lemmas.

\begin{lem}\label[lem]{fopt} Let $$f(x) =   x^2 \log\s{2 - e^{-\beta^2/x^2}} \quad \mathrm{for} \quad 0<x<1 \qquad \mathrm{and} \qquad f(0)=0. $$
Let $\alpha_1, \ldots, \alpha_k \geq  0$ with $\sum_{i=1}^k \alpha_i = 1$ then
$$ \sum_{i=1}^{k} f(\alpha_i) \leq \log 2 - \beta^2.$$	
\end{lem}	

\begin{proof} Let $F = \log 2 - \beta^2$. We have $2f(1/2)=F$.
	The function $f(x)/x$ increases on the interval $(0,1/2]$. If $\alpha_i < 1/2$ for every $i$, then $$\sum_{i=1}^{k} f(\alpha_i) \leq \sum_{i=1}^k \alpha_i \frac{f(1/2)}{1/2} = F$$ and the lemma holds.
	Otherwise, assume without loss of generality  that  $\alpha_1 \geq 1/2$. Then $$\sum_{i=2}^{k} f(\alpha_i) \leq \sum_{i=2}^k \alpha_i \frac{f(1-\alpha_1)}{1-\alpha_1} = f(1-\alpha_1),$$
	and it remains to show that $f(\alpha_1)+f(1-\alpha_1) \leq F$, i.e. to prove the lemma for $k=2$. Let $h(x)=f(x)+f(1-x).$ The function $h'(x)$  has three roots on $[0,1]$ at approximately $x_0=0.165..., 1/2,$ and $1-x_0$, and $h(x)$ decreases on $[0,x_0]$ and increases on $[x_0,1/2]$. Thus $h(x)$ is maximized at $x=0$ or $x=1/2$. We have  $h(1/2)=F=0.293\ldots$ and $h(1)=0.284\ldots < F$.      
\end{proof}

\begin{lem}\label[lem]{opt} Let $\alpha_1, \ldots, \alpha_\ell > 0$ with $\sum_{i=1}^\ell \alpha_i = 1$. Let $ 0 \leq p_1, \ldots, p_\ell < 1$. Let $ 0 \leq q_{ij}^{B}, q_{ij}^G <1$ for all $(i,j) \in [\ell]^2$. If \begin{align*} &p_i + q_{ii}^B + q_{ii}^G = 1 \qquad \mathrm{for\; all} \quad i \in [\ell], \qquad \mathrm{and} \\ &q_{ij}^B + q_{ij}^G = 1  \qquad \mathrm{for\; all} \quad i, j \in [\ell] \quad \mathrm{with} \; i \neq j,\end{align*} then at least one of the following holds:\todo{now that they're 1 and 2, maybe we should swap B and G for R and B}

\begin{itemize}
\item $\alpha_i \sqrt{-\log\s{1-p_i}} \geq \beta$ for some $i \in [\ell]$, or
\item $\sqrt{-\sum_{(i,j) \in [\ell]^2} \alpha_i\alpha_j\log\s{1-q^B_{ij}}} \geq \beta$, or
\item $\sqrt{-\sum_{(i,j)\in [\ell]^2} \alpha_i\alpha_j\log\s{1-q^G_{ij}}} \geq \beta$.
\end{itemize}
\end{lem}

\begin{proof} Suppose not, i.e. after squaring we have
\begin{equation}\label{e:first} -\alpha^2_i \log\s{1-p_i} < \beta^2 \end{equation}  for all $i \in [\ell]$,
$$-\sum_{(i,j)\in [\ell]^2} \alpha_i\alpha_j\log\s{1-q^B_{ij}}< \beta^2 \quad \mathrm{and} \quad
-\sum_{i,j \in [\ell]} \alpha_i\alpha_j\log\s{1-q^G_{ij}} < \beta^2. $$
Averaging the last two inequalities gives
$$- \sum_{i,j \in [\ell]} \alpha_i \alpha_j \frac{\log\s{1 - q_{ij}^B} - \log\s{1-q_{ij}^G}}2 < \beta^2.$$
By convexity of  $-\log\s{1-x}$ for $x < 1$, we have $$\sum_{i,j \in [\ell]^2} - \alpha_i \alpha_j \log\s{1 - \frac{q_{ij}^B + q_{ij}^G}2} \leq \sum_{i,j \in [\ell]^2} - \alpha_i \alpha_j \frac{\log\s{1 - q_{ij}^B} - \log\s{1-q_{ij}^G}}2.$$ Thus $$\sum_{i,j \in [\ell]} - \alpha_i \alpha_j \log\s{1 - \frac{q_{ij}^B + q_{ij}^G}2} < \beta^2.$$ As $q_{ii}^B + q_{ii}^G = 1 - p_i$ and $q_{ij}^B + q_{ij}^G = 1$ for $ i \neq j$ we rewrite the above $$\log 2 \cdot \sum_{i \neq j} \alpha_i \alpha_j - \sum_{i \in [k]} \alpha_i^2 \log\s{1 - \frac{1 - p_i}2} < \beta^2.$$ Meanwhile, by \eqref{e:first} we have
$$ (1 - p_i)/2 \geq e^{-\beta^2/\alpha_i^2}/2.$$
Thus
\begin{align*} \beta^2 &> \log2 \cdot \sum_{i \neq j} \alpha_i \alpha_j - \sum_{i \in [\ell]} \alpha_i^2 \log\s{1 - e^{-\beta^2/\alpha_i^2}/2} \\
&=\log 2\cdot\s{1 -\sum_{i \in [\ell]} \alpha_i^2} - \sum_{i \in [\ell]} \alpha_i^2 \log\s{1 - e^{-\beta^2/\alpha_i^2}/2} \\ &
=\log 2 - \sum_{i \in [\ell]} \alpha_i^2 \s{\log2 +\log\s{1 - e^{-\beta^2/\alpha_i^2}/2}} \\ &=
\log 2 - \sum_{i \in [\ell]} \alpha_i^2 \log\s{2 - e^{-\beta^2/\alpha_i^2}}.\end{align*}
Let
$f(x) = x^2 \log\s{2 - e^{-\beta^2/x^2}}$ be as in \cref{fopt}. Then the last inequality can be rewritten as $\sum_{i=1}^k f(\alpha_i) > \log 2 - \beta^2$, contradicting \cref{fopt}.
\end{proof}

As all the functions in \cref{opt} are continuous when the variables are bounded away from one, compactness immediately implies the following approximate version of \cref{opt}. 

\begin{cor}\label[cor]{optapprox} For every $\zeta > 0$ there exists $\delta=\delta_{\ref{optapprox}}(\zeta )$ satisfying the following. Let $\ell \leq 1/\zeta $ be a positive integer, let $\alpha_1, \ldots, \alpha_\ell > 0$ with $\sum_{i=1}^\ell \alpha_i = 1$. Let $ 0 \leq p_1, \ldots, p_\ell < 1 - \zeta $. Let $ 0 \leq q_{ij}^{B}, q_{ij}^G < 1 -\zeta $  with $q^{B}_{ij} = q^{B}_{ji},q^{G}_{ij} = q^{G}_{ji}$ for all $(i,j) \in [\ell]^2$. Then the following conditions cannot all simultaneously hold:
	\begin{itemize}
	 \item[(a)] $p_i + q_{ii}^B + q_{ii}^G  \geq  1 - \delta$ for all $i \in [\ell]$,
	 \item[(b)]  $q_{ij}^B + q_{ij}^G  \geq 1 - \delta$  for all $i, j \in [\ell]$ with $i \neq j$, 
		\item[(c)] $\alpha_i \sqrt{-\log\s{1-p_i}} \leq \beta - \zeta $ for all $i \in [\ell]$, 
		\item[(d)] $\sqrt{-\sum_{(i,j) \in [\ell]^2} \alpha_i\alpha_j\log\s{1-q^B_{ij}}} \leq \beta -\zeta $,
		\item [(e)] $\sqrt{-\sum_{(i,j)\in [\ell]^2} \alpha_i\alpha_j\log\s{1-q^G_{ij}}} \leq \beta - \zeta $.
	\end{itemize}
\end{cor}

\begin{proof}[Proof of \cref{mainthm}]
We first introduce the parameters that we will need. We assume $\eps \leq 1$ without loss of generality. Choose $\lambda < \frac{\eps \beta^2}{420 \log 2}$ then let $K = \frac{2\beta}{\lambda}$. Let $\gamma = \min\s{\gamma_{\ref{bipMinor}}(K), \gamma_{\ref{denseMinor}}(K)}$, let $\zeta = \min\s{\gamma, \frac\lambda2, \frac\beta{100}}$, and let $\xi = \delta_{\ref{optapprox}}(\zeta)$. Choose $\eta < \frac{\lambda \xi}2$ and $\eps_0 = \min\s{\lambda \xi, \lambda \eps, \zeta \beta^2}/100$. Let $\delta_0 = \delta_{\ref{twobigsets}}(\eps_0, \beta)$, let $\delta = \delta_{\ref{connectedsets}}(\lambda, \eta)$, and let $\theta = \min\s{\delta_0, \delta, \zeta}$. Finally, it will be clear that we can make $n_0$ sufficiently large based on the other parameters. %let $n_0 = \frac2\lambda \max\s{n_{0,\ref{twobigsets}}(\eps_0, \beta), n_{0,\ref{bipMinor}}(K), n_{0,\ref{denseMinor}}(K), n_{0,\ref{thomasonupper}}(\theta)}$ \textbf{todo}.
For any $n \geq n_0$, let $(G_1, G_2,G_3)$ be a 3-decomposition of $K_n$ and suppose for a contradiction that
\begin{equation}\label{e:lower2} \had(G_i) < (1- \eps) \beta \frac{n}{\sqrt{\log n}}, \end{equation}
for all $i \in [3]$. By \cref{twobigsets}  applied with parameters $\eps_0$  and $C=\beta$, we assume without loss of generality that for $j \in \{1,2\}$  \begin{itemize} \item either $|E(G_j)| \leq 3\eps_0n^2$, or \item 
there exists $P_j \subseteq V(G)$ such that $|P_j|\geq (1-\eps_0)n$ and $G_j[P_j]$ is $(\delta_0 n)$-connected for some $\delta_0 > 0$.\end{itemize} If either $P_j$ does not exist for $j \in [2]$, then let $P_j = V(G)$. Let $(V_1,V_2,\ldots,V_{k},X,Y)$ be a $(\lambda,\eta,\delta)$-controlled partition of $G_3$. If $|Y| \geq 20\lambda n$, let $V_{k+1}=Y,V_{k+2}=X$ and $\ell=k+2$, otherwise, let $V_{k+1}=X \cup Y$ and $\ell=k+1$.

We  define parameters $\alpha_i, p_i$ for $i \in [\ell]$ and  $q_{ij}^B, q_{ij}^G$ for $i, j \in [\ell]^2$. $\alpha_i$ corresponds to the relative size of $V_i$ and $p_i$ corresponds to the edge density of $V_i$ in $G_3$. Within $P_1$, $q_{ij}^B$ corresponds to the $G_1$ edge densities between $V_i$ and $V_j$, and $q_{ij}^G$ is defined analogously with $P_2$ and $G_2$. We use \eqref{e:lower2} together with \cref{thomasonupper} to show that they contradict \cref{optapprox}. 

Let $\alpha_i = \frac{|V_i|}{n}$ for $i \in [\ell]$, \begin{align*}&p_i = \frac{2|E(G_3[V_i])|}{|V_i|(|V_i|-1)}, \qquad q^{B}_{ii} = \frac{|E(G_1[V_i \cap P_1])|}{\binom{|V_i \cap P_1|}2}, \qquad q^{G}_{ii} = \frac{|E(G_2[V_i \cap P_2])|}{{\binom{|V_i \cap P_2|}2}} &\quad \mathrm{for} \; i \in [\ell-1], \; \mathrm{and} \\
 &q^{B}_{ij} = \frac{|E_{G_1}(V_i \cap P_1,V_j \cap P_1)|}{|V_i \cap P_1||V_j \cap P_1|}, \qquad  q^{G}_{ij} = \frac{|E_{G_2}(V_i \cap P_2,V_j \cap P_2)|}{|V_i \cap P_2||V_j \cap P_2|} &\quad \mathrm{for}\;    i,j \in [\ell-1], i \neq j.\end{align*}
Finally,  let $p_{\ell}=0$ and let $q^{B}_{i\ell} = q^{B}_{\ell i} = q^{G}_{i\ell} = q^{G}_{\ell i} = 1/2$ for $i \in [\ell]$. 

If $p_{i} \geq 1-\zeta$ for some $i$ then  by the choices of $\zeta$, $n_0$ and $\lambda$, \cref{denseMinor} is applicable to $G_3[V_i]$ and we have $$\had(G_3) \geq \had(G_3[V_i]) \geq K\frac{|V_i|}{\sqrt{\log |V_i|}} \geq K\lambda \frac{n}{\sqrt{\log n}} \geq \beta  \frac{n}{\sqrt{\log n}},$$
contradicting \eqref{e:lower2}. Thus $p_{i} \leq 1-\zeta$ and similarly $q^{B}_{ij},q^{G}_{ij} \leq 1-\zeta$ for $i,j \in [\ell]$, where for $i \neq j$, $i,j \in [\ell-1]$ we apply \cref{bipMinor} instead of \cref{denseMinor}. In particular, for $i \in [\ell-1]$ and $j \in [2]$, $|V_i \cap P_j| \geq (\lambda - \eps_0)n \geq \frac\lambda2 n$ by the choice of $\eps_0$, so we can repeat a similar argument to above.

It remains to show that conditions (a)--(e) of \cref{optapprox} are satisfied to obtain a contradiction.
For (a), $p_i$ is the edge density of $G_3[V_i]$. If we write $q_{ii}^{*B}$ and $q_{ii}^{*G}$ for the  edge densities of $G_1[V_i]$, $G_2[V_i]$, note that $p_i + q_{ii}^{*B} + q_{ii}^{*G} = 1$. Since $|P_j| \geq (1-\eps_0)n$ for $j \in [2]$,
\begin{align*}
q^{B}_{ii} & = \frac{|E(G_1[V_i \cap P_1])|}{\binom{|V_i \cap P_1|}2} \\
& \geq \frac{|E(G_1[V_i])|}{\binom{|V_i|}2} - \frac{ \eps_0n(|V_i|-1)}{\binom{|V_i|}2} \\
&\geq \frac{|E(G_1[V_i])|}{\binom{|V_i|}2} - \frac{2\eps_0}{\lambda}.
\end{align*}
We also have the analogue for $q_{ii}^G$. By the choice of $\eps_0$,
$$p_i + q_{ii}^B + q_{ii}^G \geq 1 - \xi.$$
For (b), we proceed similarly. If $i = \ell$ or $j = \ell$, (b) is immediate (the sum is exactly 1). Otherwise,
\begin{align*}
q^{B}_{ij} &= \frac{|E_{G_1}(V_i \cap P_1,V_j \cap P_1)|}{|V_i \cap P_1||V_j \cap P_1|} \\
& \geq \frac{|E_{G_1}(V_i,V_j)|}{|V_i||V_j|} - \frac{\eps_0 n (|V_i| + |V_j|)}{|V_i||V_j|} \\
& \geq \frac{|E_{G_1}(V_i,V_j)|}{|V_i||V_j|} - \frac{2 \eps_0}\lambda.
\end{align*}
It remains to verify that $|E_{G_3}(V_i,V_j)|$ is sufficiently small for $i, j < \ell$. Indeed, either $|E_{G_3}(V_i,V_j)| = 0$ if $Y \not\in \{V_i, V_j\}$, or
$$|E_{G_3}(V_i,Y)| \leq \eta n |Y| \Rightarrow \frac{|E_{G_3}(V_i,Y)|}{|V_i||Y|} \leq \frac{\eta}{\lambda}.$$
By the choices of $\eta$ and $\eps_0$,
$$q_{ij}^B + q_{ij}^G \geq 1 - \xi.$$
For (c), if $i = \ell$, this is immediate by $p_\ell = 0$. If $V_{\ell-1} = Y$, then $G_3[Y]$ is $(3 \lambda n)$-degenerate, so $|E(G_3[Y]) \leq 3 \lambda n |Y|$. Therefore, $$p_{\ell-1} \leq \frac{6 \lambda n}{|Y| - 1} \leq \frac{6 \lambda n}{20 \lambda n - 1} = \frac3{10} + o(1).$$
Since $\sqrt{-\log\s{1-\frac3{10}}} < 0.99\beta$,
$$\alpha_{\ell-1} \sqrt{-\log(1-p_{\ell-1})} \leq \beta - \zeta$$
by the choice of $\zeta$, for sufficiently large $n$. Finally, if $i < \ell$ and $V_i \neq Y$, then $V_i$ is $(\delta n)$-connected. By the choice of $\theta$, we can apply \cref{thomasonupper} to $G_3[V_i]$ with parameter $\theta$ and the trivial partition $(V_i)$; in particular, $p^*_{11} = p_i$. We get
$$\had(G_3) \geq \had(G_3[V_i]) \geq \frac{|V_i|}{\sqrt{\log |V_i|}} \s{\sqrt{-\log\s{1-p_i}} - \theta}.$$
By our main assumption \eqref{e:lower2} and the choices of $\zeta$ and $\theta$,
\begin{align*}
& \quad\:\: \frac{|V_i|}{\sqrt{\log |V_i|}} \s{\sqrt{-\log\s{1-p_i}} - \theta} < (1- \eps) \beta \frac{n}{\sqrt{\log n}} \\
&\Rightarrow \alpha_i \sqrt{-\log\s{1-p_i}} < (1 - \eps) \beta \sqrt\frac{\log|V_i|}{\log n} + \alpha_i \theta \\
&\Rightarrow \alpha_i \sqrt{-\log\s{1-p_i}} < (1 - \eps) \beta + \alpha_i \theta \\
&\Rightarrow \alpha_i \sqrt{-\log\s{1-p_i}} < (1 - \frac\eps2) \beta \\
&\Rightarrow \alpha_i \sqrt{-\log\s{1-p_i}} < \beta - \zeta.
\end{align*}
For (d) and (e), the argument is symmetric, so it only remains to prove (d). We can rewrite what we want to show as
$$\s{-\sum_{i, j < \ell} \alpha_i \alpha_j\log\s{1-q_{ij}^B}} + (2 \alpha_\ell - \alpha_\ell^2) \log2 \leq (\beta - \zeta)^2.$$
Since $|X| \leq \lambda n$ and $|Y| \leq 20 \lambda n$ if $Y \subseteq V_\ell$, $2\alpha_\ell \leq 42 \lambda$. Recall that $42 \lambda \log 2 \leq \frac{\eps \beta^2}{10}$. Thus, it suffices to show that
\begin{equation}\label{eq:cased}
-\sum_{i, j < \ell} \alpha_i \alpha_j \log\s{1-q_{ij}^B} < \s{1-\frac\eps2}\beta^2
\end{equation}
We first consider the case where $|E(G_1)| \leq 3 \eps_0 n^2$ and $P_1 = V(G)$. In particular,
$$\sum_{i, j < \ell} \alpha_i \alpha_j q_{ij}^B \leq \frac{2|E(G_1)|}{n^2} + \frac1n \leq 6 \eps_0 + \frac 1n \leq 7 \eps_0$$
for sufficiently large $n$. Note that, as $\frac{-\log(1-x)}x$ is increasing on $(0, 1)$ and $q_{ij}^B \leq 1 - \zeta$, we have $-\log(1-q_{ij}^B) \leq \frac{-\log(\zeta)}{1-\zeta} q_{ij}^B$ for all $i, j$. Therefore,
\begin{align*}
- \sum_{i, j < \ell} \alpha_i \alpha_j \log(1-q_{ij}^B) &\leq -\frac{\log(\zeta)}{1-\zeta} \sum_{i, j < \ell} \alpha_i \alpha_j q_{ij}^B \\
&\leq -\frac{7 \eps_0 \log(\zeta)}{1-\zeta} \\
&\leq -\frac{7\zeta\log(\zeta)}{100(1-\zeta)}\beta^2 < (1 - \frac \eps2) \beta^2
\end{align*}
by the choice of $\eps_0$, establishing \eqref{eq:cased} in this case. Now suppose that $G_1[P]$ is $(\delta_0 n)$-connected for some $|P| \geq (1-\eps_0)n$. We apply \cref{thomasonupper} to the partition $(V_1 \cap P, \ldots, V_\ell \cap P)$ of $P$, with parameter $\theta$, to deduce
$$\had(G_1[P]) \geq \frac{|P|}{\sqrt{\log |P|}} \s{\sqrt{-\sum_{i, j < \ell} \alpha_i^* \alpha_j^* \log(1-q^B_{ij})} - \theta},$$
where
$$\alpha^*_i = \frac{|V_i \cap P|}{|P|}$$
and $q_{ij}^B \leq 1 - \zeta \leq 1 - \theta$. Note that $|V_i \cap P| \geq (1-\frac{\eps_0}\lambda) |V_i|$, so
$$\had(G_1) \geq \frac{|P|}{\sqrt{\log |P|}} \s{\s{1 - \frac{\eps_0}{\lambda}} \frac{n}{|P|} \sqrt{-\sum_{i, j < \ell} \alpha_i \alpha_j \log\s{1-q^B_{ij}}} - \theta}$$
$$\Rightarrow \had(G_1) \geq \frac{n}{\sqrt{\log n}} \s{\s{1 - \frac{\eps_0}{\lambda}} \sqrt{-\sum_{i, j < \ell} \alpha_i \alpha_j \log\s{1-q^B_{ij}}} - \theta}.$$
By our main assumption \eqref{e:lower2},
$$\s{1-\frac{\eps_0}\lambda} \sqrt{-\sum_{i, j < \ell} \alpha_i \alpha_j \log(1-q_{ij}^B)} - \theta < (1-\eps) \beta$$
$$\Rightarrow \sqrt{-\sum_{i, j < \ell} \alpha_i \alpha_j \log(1-q_{ij}^B)} < \s{(1-\eps)\beta + \theta}\s{1 - \frac{\eps_0}\lambda}^{-1}$$
$$\leq \frac{1-9\eps/10}{1-\eps/10} \beta$$
$$\leq \sqrt{\s{1-\frac\eps2}}\beta,$$
proving \eqref{eq:cased}, and thus completing the proof.
%and $p^*_{ij}$ are as in \cref{thomasonupper}.

% \textbf{FIX EVERYTHING AFTER HERE}

% By \cref{thomasonupper} applied with ... to $G_3[V_i]$ for $i \in [\ell-1]$ such that $V_i \neq Y$, with a trivial partition $(V_i)$, we have $$\had(G_3[V_i]) \geq \frac{\alpha_i n}{\sqrt{\log(\alpha_i n)}}\s{\sqrt{-\log(1-p_{i})}-\eps_1}.$$
% As $\had(G_3) \geq \had(G_3[V_i])$, \eqref{e:lower2} implies $$ (1-\eps)\beta \geq \alpha_i\s{\sqrt{-\log(1-p_i)}-\eps_1}.$$\todo{here the mismatch in thm 1.5 starts to be a problem}
% As $\eps\beta > \eps_1$, we have $$ \beta > \alpha_i\sqrt{-\log(1-p_i)},$$
% implying that the first outcome of \cref{opt} cannot hold for any $i$ as above.
% We still need to show that this condition cannot hold in the case $i=\ell-1, V_i=Y$, but this is not hard. Indeed, recall that in this case $G_3[Y]$ is $(3\gamma n)$-degenerate, implying $E(G_3[Y]) \leq 3\gamma n |Y|$,  and  we also have $|Y| \geq 20\gamma n$. Thus $$p_i = \frac{|E(G_3[Y])|}{\binom{|Y|}{2}} \leq \frac{3\gamma n}{|Y|-1} \leq \frac{3\gamma n}{20\gamma n -1} \leq \frac{1}{5}.$$
% Thus $\alpha_i \sqrt{-\log(1-p_i)} \leq \sqrt{-\log\frac{4}{5}}< 1 < \beta.$

% It remains to show that the other two outcome of \cref{opt} cannot hold, i.e. by symmetry it suffices to show that 
% \begin{equation}\label{e:qbb}
% -\sum_{i,j \in [\ell]^2} \alpha_i\alpha_j\log\s{1-q^B_{ij}} < \beta^2
% \end{equation}	 
% Note that $|V_\ell| \leq |X| + 20\gamma n \leq 21\gamma n$, and so $\alpha_{\ell} \leq 21 \gamma$. As $q^B_{i\ell} = \frac{1}{2}$ for every $i \in \ell$ this implies  $$-\sum_{i \in [\ell]} \alpha_i\alpha_\ell\log\s{1-q^B_{i\ell}} \leq \alpha_{\ell}\log{2} \leq 42 \gamma \leq \frac{\eps}{2}\beta^2.$$
% Thus it further suffices to show 
% \begin{equation}\label{e:qbb2}
% 	-\sum_{i,j \in [\ell-1]^2} \alpha_i\alpha_j\log\s{1-q^B_{ij}} < \s{1 - \frac{\eps}{2}}\beta^2
% \end{equation}	
% By our assumption at the beginning of the proof either $|E(G_2)| \leq 3 \eps_0 n^2$ or there exists  $P=P_2 \subseteq [n]$ such that $|P|\geq (1-\eps_0)n$ and $G_1[P]$ is $(\delta n)$-connected. 

% Suppose first that $|E(G_1)| \leq 3 \eps_0 n^2$. Note that $\frac{-\log(1-x)}{x}$ is increasing on $(0,1)$ and $q^{B}_{ij} \leq 1-\eps_1$ for all $i,j$ we have $q^{B}_{ij} \geq \frac{\log(1-\eps_1)}{\eps_1}\log\s{1-q^{B}_{ij}}$. We have
% $$ 3 \eps_0 n^2 \geq |E(G_1)| \geq \s{\sum_{(i,j) \in [\ell-1]^2} \alpha_i\alpha_j q^{B}_{ij} }\frac{n^2}{2} - n.$$
% Also note that, as $\frac{-\log(1-x)}{x}$ is increasing on $(0,1)$ and $q^{B}_ij \leq 1-\eps_1$ for all $i,j$, we have $-\log\s{1-q^{B}_{ij}} \leq -\frac{\log(1-\eps_1)}{\eps_1}q^{B}_{ij}$. Combining all of the above we have
% \begin{align*} -\sum_{(i,j) \in [\ell-1]^2} &\alpha_i\alpha_j \log(1 - q^{B}_{ij}) \leq  -\frac{\log(1-\eps_1)}{\eps_1}\sum_{(i,j) \in [\ell-1]^2} \alpha_i\alpha_j  q^B_{ij} \\ &\leq-\frac{\log(1-\eps_1)}{\eps_1}\s{6\eps_0 +\frac{1}{n}}<  \s{1 - \frac{\eps}{2}}\beta^2, \end{align*}
% establishing \eqref{e:qbb2} in this case.

% It remains to consider the case that  $G_1[P]$ is $(\delta n)$-connected for some $|P|\geq (1-\eps_0)n$. We apply \cref{thomasonupper} to the partition $(V_1 \cap P,\ldots,V_{\ell} \cap P)$ of $P$ to deduce $$\had(G_1[P]) \geq \frac{|P|}{\sqrt{\log n}}\s{\sqrt{-\sum_{(i,j) \in [\ell-1]^2} \alpha^*_i\alpha^*_j\log\s{1-p^*_{ij}}} -\eps_2,}$$
% where  $$\alpha^*_i = \frac{|V_i \cap P|}{|P|} \leq \alpha_i \frac{n}{|P|} \leq \alpha_i (1-\eps_0)^{-1},$$ and
% $p^*_{ij}$ are as in \cref{thomasonupper}. In particular, $$1-p^{*}_{ij} = \frac{|E_{\bar{G_1}}(V_i \cap P,V_j \cap P)|}{|V_i \cap P||V_j \cap P|} \geq  \frac{|E_{\bar{G_1}}(V_i ,V_j)| - (|V_i|+|V_j|)(n-|P|)}{|V_i||V_j|} \geq (1-q^{B}_{ij}) - 2\frac{\eps_0}{\gamma} \geq  (1-q^{B}_{ij})...$$ for $i \neq j$, and similarly $p^{*}_{ii} \leq q^{B}_{ii}\s{1-\frac{\eps_0}{\gamma}}^{-2}$.

\end{proof}


%\bibliographystyle{plain}
%\bibliography{hadram}
\printbibliography


\newpage

\appendix
\counterwithin{thm}{section}
\renewcommand{\thethm}{\Alph{section}.\arabic{thm}}

\section{Proof of \cref{thomasonupper}}

\subsection{Notation}

For disjoint vertex sets $A,B\subseteq V(G)$, let $e_G(A,B)$ be the number of edges with one end in $A$ and the other in $B$, and let $e_G(A)=e(G[A])$.  We write $\kappa(G)$ for the vertex connectivity of $G$ and $\delta(G)$ for its minimum degree.

Let $\cX=(X_1,\dots,X_k)$ be a partition of $V(G)$, put
\[
 n_i=|X_i|,
 \qquad
 \alpha_i=\frac{n_i}{n},
 \qquad n=|G|,
\]
and define
\[
 p_{ii}=\frac{e_G(X_i)}{\binom{n_i}{2}},
 \qquad
 p_{ij}=\frac{e_G(X_i,X_j)}{n_i n_j}\quad(i\ne j).
\]
When $n$ is large, the hypothesis $\alpha_i\ge\eps$ ensures $n_i\ge2$, so all diagonal densities are defined.  Put
\[
 p_{ij}^*=\min\{1-\eps,p_{ij}\},
 \qquad
 q_{ij}=1-p_{ij},
 \qquad
 q_{ij}^*=1-p_{ij}^*=\max\{\eps,q_{ij}\},
\]
and
\begin{equation}\label{eq:Lambda-star-definition}
 \Lambda^*(G,\cX)
 :=-\sum_{i,j=1}^k \alpha_i\alpha_j\log q_{ij}^*.
\end{equation}
The sum is over ordered pairs.  Thus, because $p_{ij}=p_{ji}$, every off-diagonal unordered pair contributes twice.

A partition $W_1,\dots,W_t$ of $V(G)$ is \emph{complete} if there is an edge between $W_a$ and $W_b$ whenever $a\ne b$, and it is an \emph{equipartition} if its part sizes differ by at most one.

The proof follows the architecture of Thomason's complete-equipartition and dense-minor arguments \cite[Theorem 3.1, Corollary 3.1, and Theorem 4.1]{MR1814910}.

\subsection{A multitype complete-equipartition theorem}

The next theorem is a multitype analogue of Thomason's complete-equipartition theorem.  The parameters $\ell_i$ specify how many vertices each random part receives from $X_i$. \todo{the analogue in the Thomason paper is actually a theorem, but changed it to lemma for now}

\begin{lem}\label[lem]{lem:multitype-finite}
Let $G$ be a graph whose vertex set is partitioned as $X_1\cup\cdots\cup X_k$.  Suppose that
\[
 |X_i|=s\ell_i\qquad(1\le i\le k),
\]
where $s\ge2$ and $\ell_i\ge1$ are integers, and put $\ell=\sum_i\ell_i$.  Let $q_{ij}=1-p_{ij}$ be the complementary densities of $G$ with respect to this partition.  Fix parameters
\[
 0<\tau<1,
 \qquad
 0<\varphi<1,
 \qquad
 \omega\ge1,
\]
and set
\[
 \overline q_{ij}:=\min\{1,q_{ij}+\tau\}
\]
and
\begin{equation}\label{eq:P-def}
 P:=(2\omega)^\ell\,
 \tau^{-\ell(k+\varphi\ell)}
 \prod_{i,j=1}^k \overline q_{ij}^{\ell_i\ell_j}.
\end{equation}
Then $G$ has a complete equipartition into at least
\begin{equation}\label{eq:finite-number-parts}
 s-\frac{4s}{\omega\varphi}-4s^2P
\end{equation}
parts and at most $s$ parts.
\end{lem}

\begin{proof}
If the quantity in \eqref{eq:finite-number-parts} is at most one, the trivial one-part equipartition proves the assertion.  We may therefore assume that it is larger than one.

For $u\in X_i$ and $1\le j\le k$, let
\[
 Q_j(u):=\{v\in X_j\setminus\{u\}: uv\notin E(G)\},
 \qquad
 d_j(u):=\frac{|Q_j(u)|}{|X_j|},
\]
and put
\[
 \overline d_j(u):=\min\{1,d_j(u)+\tau\},
 \qquad
 \rho(u):=\prod_{j=1}^k \overline d_j(u)^{\ell_j/\ell}.
\]
Thus $\tau\le\rho(u)\le1$ for every vertex $u$.

Fix $i$.  If $j\ne i$, then
\[
 \frac1{|X_i|}\sum_{u\in X_i}d_j(u)=q_{ij}.
\]
For $j=i$, the denominator in $d_i(u)$ is $|X_i|$, whereas $q_{ii}$ is normalized by $\binom{|X_i|}{2}$; hence
\[
 \frac1{|X_i|}\sum_{u\in X_i}d_i(u)
 =q_{ii}\frac{|X_i|-1}{|X_i|}\le q_{ii}.
\]
Since $\overline d_j(u)\le d_j(u)+\tau$ and $\overline d_j(u)\le1$, it follows that
\[
 \frac1{|X_i|}\sum_{u\in X_i}\overline d_j(u)\le \overline q_{ij}.
\]
The geometric-mean--arithmetic-mean inequality therefore gives
\begin{align}
 \left(\prod_{u\in X_i}\rho(u)\right)^{1/|X_i|}
 &=\prod_{j=1}^k
 \left(\prod_{u\in X_i}\overline d_j(u)\right)^{\ell_j/(\ell|X_i|)}
 \notag\\
 &\le \prod_{j=1}^k \overline q_{ij}^{\ell_j/\ell}.
 \label{eq:rho-geomean}
\end{align}

Order the vertices in each class $X_i$ so that $\rho$ is nondecreasing, and divide $X_i$ consecutively into blocks
\[
 B_{i,1},\dots,B_{i,\ell_i},
 \qquad |B_{i,a}|=s.
\]
Choose, independently for every block, a uniformly random ordering of its vertices.  For $1\le t\le s$, let $W_t$ consist of the $t$th vertex in every block.  Thus $W_1,\dots,W_s$ is a random partition into $s$ sets of size $\ell$, each set containing exactly $\ell_i$ vertices from $X_i$.

For $W\subseteq V(G)$, define
\[
 N(W):=\{u\in V(G): W\subseteq Q(u)\},
 \qquad Q(u):=\bigcup_{j=1}^k Q_j(u).
\]
Thus $u\in N(W)$ precisely when $u\notin W$ and $u$ has no neighbor in $W$.  For a fixed vertex $u$ and a fixed random part $W_t$, independence between blocks and the arithmetic-geometric mean inequality within the blocks belonging to each $X_j$ yield
\begin{align*}
 \Pr(W_t\subseteq Q(u))
 &=\prod_{j=1}^k\prod_{a=1}^{\ell_j}
 \frac{|Q_j(u)\cap B_{j,a}|}{s}\\
 &\le\prod_{j=1}^k
 \left(\frac{|Q_j(u)|}{s\ell_j}\right)^{\ell_j}
 =\prod_{j=1}^k d_j(u)^{\ell_j}
 \le \rho(u)^\ell.
\end{align*}
Let
\[
 \theta_{i,a}:=\max_{u\in B_{i,a}}\rho(u).
\]
It follows that, for every random part $W_t$,
\begin{equation}\label{eq:block-expectation}
 \mathbb E\,|B_{i,a}\cap N(W_t)|\le s\theta_{i,a}^{\ell}.
\end{equation}

Call a block $B_{i,a}$ \emph{eligible} if $a<\ell_i$.  We say that $W_t$ rejects an eligible block $B_{i,a}$ when
\[
 |B_{i,a}\cap N(W_t)|>\omega s\theta_{i,a}^{\ell}.
\]
By \eqref{eq:block-expectation} and Markov's inequality, each eligible block is rejected with probability at most $1/\omega$.  Call $W_t$ \emph{acceptable} if it rejects at most $\varphi\ell$ eligible blocks.  Since there are at most $\ell$ eligible blocks,
\begin{equation}\label{eq:unacceptable-prob}
 \Pr(W_t\text{ is unacceptable})\le\frac1{\omega\varphi}.
\end{equation}

Fix an acceptable part $W=W_t$.  For every $i$, let $M_i$ be the set of indices $a<\ell_i$ for which $B_{i,a}$ is not rejected, put $m_i=|M_i|$, and let $m=\sum_i m_i$.  Acceptability gives
\begin{equation}\label{eq:m-lower}
 m\ge \ell-k-\varphi\ell,
 \qquad\text{and hence}\qquad
 \ell-m\le k+\varphi\ell.
\end{equation}
For $a<\ell_i$, every vertex of $B_{i,a+1}$ has $\rho$-value at least $\theta_{i,a}$.  Consequently,
\[
 \theta_{i,a}^{s}\le\prod_{u\in B_{i,a+1}}\rho(u).
\]
Multiplying over $a\in M_i$ and using $\rho(u)\ge\tau$ on omitted blocks gives
\begin{align*}
 \left(\prod_{a\in M_i}\theta_{i,a}\right)^s
 &\le \prod_{a\in M_i}\prod_{u\in B_{i,a+1}}\rho(u)\\
 &\le \left(\prod_{u\in X_i}\rho(u)\right)
 \tau^{-s(\ell_i-m_i)}.
\end{align*}
Taking $s$th roots and invoking \eqref{eq:rho-geomean}, we obtain
\begin{equation}\label{eq:block-max-product-i}
 \prod_{a\in M_i}\theta_{i,a}
 \le
 \left(\prod_{j=1}^k\overline q_{ij}^{\ell_j/\ell}\right)^{\ell_i}
 \tau^{-(\ell_i-m_i)}.
\end{equation}
Multiplying \eqref{eq:block-max-product-i} over $i$ and then raising to the power $\ell$ yields
\begin{equation}\label{eq:block-max-product-all}
 \prod_{i=1}^k\prod_{a\in M_i}\theta_{i,a}^{\ell}
 \le
 \left(\prod_{i,j=1}^k\overline q_{ij}^{\ell_i\ell_j}\right)
 \tau^{-\ell(\ell-m)}.
\end{equation}

Let $W_{t'}$ be a second part, where $t'\ne t$.  Conditional on $W_t=W$, the chosen vertex of $W_{t'}$ in each block is uniform among the $s-1$ vertices not used by $W$, independently from block to block.  Also $W\cap N(W)=\varnothing$, because $u\notin Q(u)$ for every vertex $u$.  Therefore, using only the nonrejected eligible blocks,
\begin{align*}
 \Pr(\text{no edge exists between $W_{t'}$ and $W$})
 &=\Pr(W_{t'}\subseteq N(W))\\
 &\le\prod_{i=1}^k\prod_{a\in M_i}
 \frac{|B_{i,a}\cap N(W)|}{s-1}\\
 &\le (2\omega)^\ell
 \prod_{i=1}^k\prod_{a\in M_i}\theta_{i,a}^{\ell}\\
 &\le P,
\end{align*}
where the last inequality follows from \eqref{eq:m-lower}, \eqref{eq:block-max-product-all}, and the definition \eqref{eq:P-def}.

Let $U$ be the number of unacceptable parts.  By \eqref{eq:unacceptable-prob},
\[
 \mathbb E U\le \frac{s}{\omega\varphi}.
\]
Let $Z$ be the number of ordered pairs $(W_{t'},W_t)$ of acceptable parts for which there is no edge between $W_{t'}$ and $W_t$.  The preceding conditional estimate gives
\[
 \mathbb E Z\le s(s-1)P<s^2P.
\]
By Markov's inequality and the union bound, there is an outcome for which simultaneously
\[
 U\le\frac{4s}{\omega\varphi}
 \qquad\text{and}\qquad
 Z\le4s^2P.
\]
Delete all unacceptable parts.  Then, for each remaining defective ordered pair, delete one of its two parts.  At most $Z$ further parts are deleted, and the surviving parts form a complete partition.  Finally distribute all vertices in deleted parts among the surviving parts so that their sizes differ by at most one.  Adding vertices cannot destroy an edge already present between two surviving parts.  This gives a complete equipartition with at least the number of parts in \eqref{eq:finite-number-parts}.
\end{proof}

We shall use the following elementary stability estimate twice: first to pass to a subgraph whose class sizes are exact multiples of prescribed integers, and later after the two reservoirs have been removed.

\begin{lem}\label[lem]{lem:stability}
For every $a>0$ there is a constant $C_a$ with the following property.  Let $G$ be a graph of order $n$, with a partition $X_1,\dots,X_k$ satisfying $|X_i|\ge an$ for every $i$.  Delete a set $Z$ of at most $\delta n$ vertices, where $0\le\delta<a/2$, and put
\[
 G'=G-Z,
 \qquad
 X_i'=X_i\setminus Z,
 \qquad
 \alpha_i'=\frac{|X_i'|}{|G'|}.
\]
Let $p_{ij}',q_{ij}'$ be the induced densities in $G'$ with respect to $X_1',\dots,X_k'$.  Then, for all sufficiently large $n$,
\[
 \max_i|\alpha_i'-\alpha_i|\le C_a\delta,
 \qquad
 \max_{i,j}|p_{ij}'-p_{ij}|=
 \max_{i,j}|q_{ij}'-q_{ij}|\le C_a\delta.
\]
Consequently, if $q_{ij}\ge a$ for all $i,j$ and $\delta$ is sufficiently small in terms of $a$, then $q_{ij}'\ge a/2$ and
\begin{equation}\label{eq:Lambda-stability}
 \left|
 -\sum_{i,j}\alpha_i'\alpha_j'\log q_{ij}'
 +\sum_{i,j}\alpha_i\alpha_j\log q_{ij}
 \right|
 \le C_a\delta.
\end{equation}
\end{lem}

\begin{proof}
We will choose $C_a$ sufficiently large to satisfy the inequalities throughout the proof. It will be clear that such a choice is possible and independent of $G$ and other parameters, except $a$. 

Write $n_i=|X_i|$, $z_i=|Z\cap X_i|$, $n'=n-|Z|$, and $n_i'=n_i-z_i$. Note that $n_i'\ge(a-\delta)n\ge an/2$.

For $i\ne j$, write $A=e_G(X_i,X_j)$, $A'=e_{G'}(X_i',X_j')$, $R=n_i n_j$, and $R'=n_i'n_j'$.  Then
\[
 0\le A-A'\le z_i n_j+z_j n_i,
 \qquad
 0\le R-R'\le z_i n_j+z_j n_i.
\]
Since $A\le R$,
\begin{align*}
 |p_{ij}'-p_{ij}|
 &\le \frac{|A-A'|}{R'}+A\left|\frac1{R'}-\frac1R\right|\\
 &\le \frac{2(z_i n_j+z_j n_i)}{n_i'n_j'}
 \le \frac{16}{a^2}\delta.
\end{align*}
For $i=j$, put $R=\binom{n_i}{2}$ and $R'=\binom{n_i'}{2}$.  Both the number of deleted possible pairs and $R-R'$ are at most $z_i n_i$.  The same calculation, with a changed constant, gives $|p_{ii}'-p_{ii}|\le C_a\delta$ for all sufficiently large $n$.

Also
\[
 \left|\frac{n_i-z_i}{n-|Z|}-\frac{n_i}{n}\right|
 =\frac{|n_i|Z|-nz_i|}{n(n-|Z|)}\le \frac{\delta}{1-a/2},
\]
which proves the first two assertions. If $q_{ij}\ge a$ and $C_a\delta\le a/2$, then $q_{ij}'\ge a/2$.  On $[a/2,1]$ the function $x\mapsto-\log x$ is $2/a$-Lipschitz and is bounded by $\log(2/a)$.  Moreover $k\le1/a$, so
\[
 \sum_{i,j}|\alpha_i'\alpha_j'-\alpha_i\alpha_j|
 \le 2\sum_i|\alpha_i'-\alpha_i|
 \le \frac2a \max_i|\alpha_i'-\alpha_i|
\]
Therefore,
\begin{align*}
\left|
 -\sum_{i,j}\alpha_i'\alpha_j'\log q_{ij}'
 +\sum_{i,j}\alpha_i\alpha_j\log q_{ij}
 \right|
 &\leq\left|\sum_{i,j} \alpha_i'\alpha_j'(\log q_{ij}-\log q'_{ij})\right| \\
 &+\left|\sum_{i,j} (\alpha_i\alpha_j - \alpha_i'\alpha_j') \log q_{ij}\right| \\
 &\leq \frac2a\max_{i,j}\left|q'_{ij}-q_{ij}\right|\\
 &+ \log\s{\frac2a}\s{\frac2a\max_i\left|\alpha_i'-\alpha_i\right|}.
\end{align*}
Thus, enlarging $C_a$ proves \eqref{eq:Lambda-stability}.
\end{proof}

\begin{cor}\label[cor]{cor:multitype-asymptotic}
For every $a>0$, $\lambda>0$, and $\eta>0$, there are constants $N=N(a,\lambda,\eta)$ and $C=C(a,\lambda,\eta)$ with the following property.  Let $G$ be a graph of order $n\ge N$, and let $X_1,\dots,X_k$ be a partition of $V(G)$.  Suppose that
\[
 \alpha_i\ge a,
 \qquad
 q_{ij}\ge a
 \qquad(1\le i,j\le k),
\]
and put
\[
 \Lambda=-\sum_{i,j=1}^k\alpha_i\alpha_j\log q_{ij}.
\]
If $\Lambda\ge\lambda$, then $G$ has a complete equipartition into $t$ parts, where
\begin{equation}\label{eq:cor-lower}
 t\ge(1-\eta)n\sqrt{\frac{\Lambda}{\log n}},
 \qquad
 t\le C\frac{n}{\sqrt{\log n}}.
\end{equation}
In particular, every part has at most $C\sqrt{\log n}$ vertices, after increasing $C$ if necessary.
\end{cor}

\begin{proof}
Because $\alpha_i\ge a$, we have $k\le1/a$, and because $q_{ij}\ge a$,
\begin{equation}\label{eq:Lambda-upper}
 \lambda\le\Lambda\le\log(1/a).
\end{equation}
Choose constants $\gamma,\xi>0$ so small that
\begin{equation}\label{eq:gamma-delta-choices}
 \frac{1-2\xi}{1+\gamma}>1-\eta,
\end{equation}
\begin{equation}\label{eq:gamma-delta-choices2}
(1+\gamma)^2(1-3\xi)>1+2\xi.
\end{equation} Choose $\tau\in(0,1)$ so small that $\tau/a\le\xi\lambda$, then choose $\varphi>0$ so small that
\[
 \varphi\log(1/\tau)\le\xi\lambda,
\]
and finally choose $\omega\ge1$ so large that
\begin{equation}\label{eq:omega-choice}
 \frac4{\omega\varphi}\le\xi.
\end{equation}
All these choices depend only on $a,\lambda,\eta$. Set
\begin{equation}\label{eq:ell-choice}
 \ell=\left\lceil(1+\gamma)\sqrt{\frac{\log n}{\Lambda}}\right\rceil.
\end{equation}
Choose positive integers $\ell_1,\dots,\ell_k$ with $\sum_i\ell_i=\ell$ and
\begin{equation}\label{eq:ell-rounding}
 \left|\frac{\ell_i}{\ell}-\alpha_i\right|\le\frac1\ell.
\end{equation}
Such a choice is obtained by rounding the numbers $\alpha_i\ell$ and distributing the rounding remainder; positivity holds for large $n$ because $\alpha_i\ge a$ and $\ell\to\infty$.   Let
\[
 s:=\min_i\left\lfloor\frac{|X_i|}{\ell_i}\right\rfloor,
\]
and choose arbitrary subsets $X_i^0\subseteq X_i$ of size $s\ell_i$.  Put $G_0=G[\bigcup_iX_i^0]$.  From \eqref{eq:ell-rounding}, $|\ell_i-\alpha_i\ell|\le1$.  Since $\alpha_i\ge a$, this implies, uniformly in $i$,
\[
 \frac{|X_i|}{\ell_i}=\frac{\alpha_i n}{\alpha_i\ell+O(1)}
 =\frac n\ell\left(1+O_a(1/\ell)\right).
\]
Taking the minimum and accounting for the floor changes the result by at most one, so
\begin{equation}\label{eq:s-asymptotic}
 s=\frac n\ell\left(1+O_a(1/\ell)\right),
 \qquad
 |G_0|=s\ell=n\left(1+O_a(1/\ell)\right).
\end{equation}
Thus only $O_a(n/\ell)$ vertices were deleted; if we say that $\delta n$ vertices were deleted, then $\delta$ can be made arbitrarily small with large enough $\ell$. Let $q_{ij}^0$ be the complementary densities in $G_0$, define
\[
 \Lambda_0^{\tau}
 :=-\sum_{i,j=1}^k\frac{\ell_i\ell_j}{\ell^2}
 \log\min\{1,q_{ij}^0+\tau\}.
\]
Define $g_\tau(x)=-\log\min\{1,x+\tau\}$.  For $x\in[a,1]$,
\[
 0\le -\log x-g_\tau(x)\le\frac\tau a.
\]
Thus, with $\Lambda^\tau=\sum_{i,j}\alpha_i\alpha_j g_\tau(q_{ij})$,
\[
 \Lambda^\tau\ge\Lambda-\frac\tau a\ge(1-\xi)\Lambda.
\]
Lemma~\ref{lem:stability} applied to $G$ and $G_0$ yields $\max_{i,j}\left|q_{ij}^0 - q_{ij}\right| = o_a(1)$ and $q^0_{ij} \geq a/2$. The function $g_\tau$ is $2/a$-Lipschitz on $[a/2,1]$, so \eqref{eq:ell-rounding} and $k\le1/a$ imply $|\Lambda_0^\tau-\Lambda^\tau|=o(1)$ uniformly over all admissible inputs. Since $\Lambda\ge\lambda$, for all sufficiently large $n$ this error is at most $\xi\lambda\le\xi\Lambda$.  Hence
\begin{equation}\label{eq:regularized-Lambda-lower}
 \Lambda_0^{\tau}\ge(1-2\xi)\Lambda.
\end{equation}

Apply Lemma~\ref{lem:multitype-finite} to $G_0$, with the integers $s,\ell_i$ and the fixed parameters $\tau,\varphi,\omega$.  For the quantity $P$ in \eqref{eq:P-def}, \eqref{eq:regularized-Lambda-lower} and the choice of $\varphi$ give
\begin{align*}
 \log P
 &=\ell\log(2\omega)
   +\ell(k+\varphi\ell)\log(1/\tau)
   -\ell^2\Lambda_0^{\tau}\\
 &\le -(1-3\xi)\Lambda\ell^2+O_{a,\lambda,\eta}(\ell).
\end{align*}
By \eqref{eq:ell-choice}, \eqref{eq:gamma-delta-choices2}, and $\ell=O_\lambda(\sqrt{\log n})$, it follows that, for large $n$,
\begin{equation}\label{eq:P-small}
 P\le n^{-1-\xi}.
\end{equation}
Lemma~\ref{lem:multitype-finite}, \eqref{eq:omega-choice}, and \eqref{eq:P-small} therefore give a complete equipartition of $G_0$ into at least
\[
 s\left(1-\xi-4sP\right)\ge(1-2\xi)s
\]
parts.  Add the vertices of $G-G_0$ to these parts, always adding the next vertex to a currently smallest part.  Completeness is preserved and the final partition is an equipartition of $V(G)$.

Using \eqref{eq:s-asymptotic}, \eqref{eq:ell-choice}, and \eqref{eq:gamma-delta-choices}, the number $t$ of parts satisfies
\[
 t\ge(1-\eta)n\sqrt{\frac\Lambda{\log n}}
\]
for all sufficiently large $n$.  On the other hand, the construction has $t\le s\le n/\ell+1$, and \eqref{eq:Lambda-upper} gives
\[
 t\le C\frac n{\sqrt{\log n}}
\]
for a constant $C=C(a,\lambda,\eta)$.  The lower bound with $\Lambda\ge\lambda$ gives $n/t=O_{\lambda,\eta}(\sqrt{\log n})$, proving the asserted part-size bound.
\end{proof}

\subsection{Reservoir lemmas}

We shall use the following standard Chernoff estimates: if $Y$ is a sum of independent Bernoulli random variables with mean $\mu$, then
\begin{equation}\label{eq:Chernoff}
 \Pr(Y\le\mu/2)\le e^{-\mu/8},
 \qquad
 \Pr(Y\ge2\mu)\le e^{-\mu/3}.
\end{equation}
These inequalities are standard consequences of the exponential-moment method; see, for example, Theorem 4.4 from \cite{Mitzenmacher_Upfal_2005}.

\begin{lem}\label[lem]{lem:short-paths}
Let $G$ be an $n$-vertex graph with vertex connectivity $\kappa$.  Every two distinct vertices $u,v$ are joined by at least $\kappa/2$ internally vertex-disjoint paths of length at most
\[
 \frac{2n}{\kappa}+1.
\]
\end{lem}

\begin{proof}
By Menger's theorem, there are $\kappa$ internally vertex-disjoint $u$--$v$ paths; see, for example, \cite[Section 3.3]{MR3644391}.  Their sets of internal vertices are pairwise disjoint and together contain at most $n-2$ vertices.  If more than $\kappa/2$ of the paths had more than $2n/\kappa$ internal vertices, their internal vertices would total more than $n$, a contradiction.  Thus at least $\kappa/2$ paths have at most $2n/\kappa$ internal vertices, and hence length at most $2n/\kappa+1$.
\end{proof}

\begin{lem}\label[lem]{lem:reservoirs}
Fix $c>0$, and put
\[
 H:=\left\lceil\frac2c\right\rceil+1,
 \qquad
 r:=\frac1{\log\log n}.
\]
For all sufficiently large $n$, every $n$-vertex graph $G$ with $\kappa(G)\ge cn$ has disjoint sets $C,D\subseteq V(G)$ such that
\begin{enumerate}
 \item $|C|,|D|\le2rn$;
 \item every vertex $v$ has at least $r\deg_G(v)/2$, and hence at least $crn/2$, neighbors in $D$;
 \item for every two distinct vertices $u,v$, there are at least
 \begin{equation}\label{eq:T-def}
  T:=\frac{c}{4}nr^{H-1}
 \end{equation}
 internally vertex-disjoint $u$--$v$ paths of length at most $H$ whose internal vertices all lie in $C$.
\end{enumerate}
\end{lem}

\begin{proof}
For each vertex, independently assign one of the labels $C,D,R$ with probabilities $r,r,1-2r$, respectively.  The first assertion fails with probability at most $2e^{-rn/3}$ by \eqref{eq:Chernoff}.

For a fixed vertex $v$, the number of its neighbors labeled $D$ is binomial with mean $r\deg_G(v)\ge crn$, since $\delta(G)\ge\kappa(G)\ge cn$.  Thus the probability that it is smaller than $r\deg_G(v)/2$ is at most $e^{-crn/8}$.  A union bound over all vertices shows that the second assertion fails with probability at most $ne^{-crn/8}$.

Fix distinct $u,v$.  By Lemma~\ref{lem:short-paths}, there is a family of at least $cn/2$ internally disjoint $u$--$v$ paths, each of length at most $H$.  For each path, the probability that all of its internal vertices receive label $C$ is at least $r^{H-1}$.  These events are independent because the paths have disjoint sets of internal vertices.  Their number has mean at least $cnr^{H-1}/2$, so by \eqref{eq:Chernoff} the probability that fewer than $T$ survive is at most
\[
 \exp\left(-\frac{cnr^{H-1}}{16}\right).
\]
A union bound over fewer than $n^2$ pairs completes the proof, because
\[
 nr^{H-1}=\frac{n}{(\log\log n)^{H-1}}\gg\log n.
\]
For large $n$, the sum of all three failure probabilities is less than one.
\end{proof}

\begin{lem}\label[lem]{lem:dominating-set}
Let $F$ be a bipartite graph with bipartition $(A,B)$.  Suppose every vertex of $A$ has at least $\theta|B|$ neighbors in $B$, where $0<\theta<1$.  Then there is a set $M\subseteq B$ meeting the neighborhood of every vertex of $A$ and satisfying
\begin{equation}\label{eq:dominating-bound}
 |M|\le \left\lceil\frac{\log |A|}{-\log(1-\theta)}\right\rceil+1.
\end{equation}
\end{lem}

\begin{proof}
Starting with $A_0=A$, choose $b_{t+1}\in B$ to maximize the number of neighbors it has in $A_t$, and let $A_{t+1}$ be the set of vertices of $A_t$ not adjacent to $b_{t+1}$.  There are at least $\theta|A_t||B|$ edges between $A_t$ and $B$, so some vertex of $B$ has at least $\theta|A_t|$ neighbors in $A_t$.  Hence
\[
 |A_{t+1}|\le(1-\theta)|A_t|,
 \qquad
 |A_t|\le(1-\theta)^t|A|.
\]
The right side is smaller than one once $t$ exceeds the quantity in \eqref{eq:dominating-bound} minus one.  The selected vertices then form the desired set $M$.
\end{proof}

\begin{proof}[Proof of Theorem~\ref{thomasonupper}]
For $\eps\ge1$ the hypotheses are empty for all $n\ge2$, since $\kappa(G)\le n-1$.  Hence assume $0<\eps<1$.  Let
\[
 \Lambda^*=-\sum_{i,j=1}^k\alpha_i\alpha_j\log q_{ij}^*.
\]
Because $\alpha_i\ge\eps$, we have $k\le1/\eps$, and because $q_{ij}^*\ge\eps$,
\begin{equation}\label{eq:Lambda-main-upper}
 \Lambda^*\le\log(1/\eps).
\end{equation}
We also need a lower bound depending only on $\eps$.  Since $-\log(1-x)\ge x$ for $0\le x<1$ and $p_{ij}^*\ge(1-\eps)p_{ij}$ for all $i,j$,
\[
 \Lambda^*\ge\sum_{i,j=1}^k\alpha_i\alpha_jp_{ij}^*
 \ge(1-\eps)\sum_{i,j=1}^k\alpha_i\alpha_jp_{ij}.
\]
The ordered sum on the right counts each off-diagonal edge twice, and for a diagonal class
\[
 \alpha_i^2p_{ii}
 =\frac{n_i^2}{n^2}\frac{e_G(X_i)}{\binom{n_i}{2}}
 =\frac{2e_G(X_i)}{n^2}\frac{n_i}{n_i-1}
 \ge\frac{2e_G(X_i)}{n^2}.
\]
Consequently
\begin{equation}\label{eq:Lambda-main-lower}
 \Lambda^*\ge (1-\eps)\frac{2e(G)}{n^2}
 \ge (1-\eps)\frac{\delta(G)}n \ge (1-\eps)\frac{\kappa(G)}n
 \ge \eps(1-\eps).
\end{equation}
For $i<j$, choose exactly
\[
 \left\lfloor p_{ij}^* n_i n_j\right\rfloor
\]
edges from $G$ between $X_i$ and $X_j$.  For $i=j$, choose exactly
\[
 \left\lfloor p_{ii}^*\binom{n_i}{2}\right\rfloor
\]
edges from $G[X_i]$.  This is possible because $p_{ij}^*\le p_{ij}$.  Let $H$ be the spanning subgraph consisting of all chosen edges.  Denote the densities of $H$ by $\widehat p_{ij}$, put $\widehat q_{ij}=1-\widehat p_{ij}$, and define
\[
 \widehat\Lambda=-\sum_{i,j=1}^k\alpha_i\alpha_j\log \widehat q_{ij}.
\]
Then $\widehat p_{ij}\le p_{ij}^*\le1-\eps$, so $\widehat q_{ij}\ge\eps$.  Moreover, since every denominator is at least a constant multiple of $n^2$ and $k\le1/\eps$,
\begin{equation}\label{eq:hatLambda-close}
 \widehat\Lambda=\Lambda^*+o_{\eps}(1),
 \qquad
 \frac{\widehat\Lambda}{\Lambda^*}=1-o_{\eps}(1),
\end{equation}
where the ratio statement uses \eqref{eq:Lambda-main-lower}.  Note that the final minor will be found in $G$; $H$ is used only to obtain a complete equipartition with bounded complementary densities.

Apply Lemma~\ref{lem:reservoirs} to the original graph $G$ with $c=\eps$, obtaining disjoint sets $C,D$ and constants
\[
 H_0=\left\lceil\frac2\eps\right\rceil+1,
 \qquad
 r=\frac1{\log\log n},
 \qquad
 T=\frac\eps4nr^{H_0-1}.
\]
Let
\[
 H':=H-C-D,
 \qquad
 X_i':=X_i\setminus(C\cup D),
 \qquad
 n':=|H'|.
\]
Since $|C|+|D|\le4rn=o(n)$, Lemma~\ref{lem:stability} applied to $H$ and $H'$ with $a=\eps$, \eqref{eq:Lambda-main-lower}, and the inequalities $\alpha_i\ge\eps$ and $\widehat q_{ij}\ge\eps$ imply, uniformly over all admissible graphs and partitions,
\[
\begin{aligned}
 \alpha_i'\ge\eps/2,
 \qquad
 q_{ij}'\ge\eps/2,
 \qquad
 \Lambda':=-\sum_{i,j}\alpha_i'\alpha_j'\log q_{ij}'\ge\frac{\eps(1-\eps)}4,
\end{aligned}
\]
\begin{equation}\label{eq:prime-stability}
\frac{\Lambda'}{\Lambda^*}=1-o_{\eps}(1),
\end{equation}
where $q_{ij}'$ denotes the complementary density in $H'$ with respect to $X_i',X_j'$ and $\alpha_i'=|X_i'|/n'$.  Also $n'/n=1-o(1)$.

Apply Corollary~\ref{cor:multitype-asymptotic} to $H'$ with
\[
 a=\eps/2,
 \qquad
 \lambda=\eps(1-\eps)/4,
 \qquad
 \eta=\eps/4.
\]
We obtain a complete equipartition
\[
 V(H')=W_1\cup\cdots\cup W_m
\]
with
\begin{equation}\label{eq:m-partition-bounds}
 m\ge(1-\eps/4)n'\sqrt{\frac{\Lambda'}{\log n'}},
\end{equation}
\begin{equation}\label{eq:m-partition-bounds2}
 m\le C_\eps\frac n{\sqrt{\log n}},
\end{equation}
\begin{equation}\label{eq:m-partition-bounds3}
 \max_{a \in [m]} |W_a|\le C_\eps\sqrt{\log n}.
\end{equation}
By \eqref{eq:prime-stability}, after increasing $n_0$ we may ensure that
\[
 \frac{n'}n
 \sqrt{\frac{\Lambda'}{\Lambda^*}}
 \sqrt{\frac{\log n}{\log n'}}
 \ge1-\frac\eps2.
\]
Together with \eqref{eq:m-partition-bounds}, this gives
\begin{equation}\label{eq:m-target}
 m\ge(1-\eps)\frac n{\sqrt{\log n}}\sqrt{\Lambda^*},
\end{equation}
because $(1-\eps/4)(1-\eps/2)\ge1-\eps$ for $0<\eps<1$.
It remains to turn the complete partition of the trimmed graph $H'$ into connected branch sets in the original graph $G$.

Put
\begin{equation}\label{eq:L-def}
 L:=\left\lceil
 \frac{\log(C_\eps\sqrt{\log n})}{-\log(1-\eps/8)}
 \right\rceil+1.
\end{equation}
Thus $L=O_\eps(\log\log n)$.  By \eqref{eq:m-partition-bounds2},
\begin{equation}\label{eq:ML-small}
 mL=O_\eps\left(\frac{n\log\log n}{\sqrt{\log n}}\right)=o_\eps(rn).
\end{equation}
In particular, for large $n$,
\begin{equation}\label{eq:ML-less-D}
 mL\le\frac{\eps rn}{4}.
\end{equation}

We construct pairwise disjoint sets $M_1,\dots,M_m\subseteq D$, one at a time, so that $M_a$ dominates $W_a$ in $G$ and $|M_a|\le L$.  Suppose $M_1,\dots,M_{a-1}$ have been chosen and put $U=\bigcup_{b<a}M_b$.  Lemma~\ref{lem:reservoirs}(2) and \eqref{eq:ML-less-D} imply that every vertex of $W_a$ has at least
\[
 \frac{\eps rn}{2}-|U|\ge\frac{\eps rn}{4}
\]
neighbors in $D\setminus U$.  Since $|D\setminus U|\le|D|\le2rn$, this is at least an $\eps/8$ fraction of $D\setminus U$.  Lemma~\ref{lem:dominating-set}, \eqref{eq:m-partition-bounds3}, and \eqref{eq:L-def} now give a set
\[
 M_a\subseteq D\setminus U,
 \qquad |M_a|\le L,
\]
which meets the neighborhood in $G$ of every vertex of $W_a$.  This completes the induction.

For every $a$, list $M_a=\{x_{a,1},\dots,x_{a,t_a}\}$.  We connect consecutive vertices $x_{a,j},x_{a,j+1}$ by paths whose internal vertices lie in $C$, choosing all internal vertices globally disjoint.  The total number of requested paths is at most $mL$, and each uses at most $H_0-1$ internal vertices.  By \eqref{eq:ML-small} and the definition of $T$,
\begin{equation}\label{eq:used-C-less-T}
 (H_0-1)mL
 =O_\eps\left(\frac{n\log\log n}{\sqrt{\log n}}\right)
 <\frac\eps4\frac{n}{(\log\log n)^{H_0-1}}
 =T
\end{equation}
for all sufficiently large $n$.
At any stage, Lemma~\ref{lem:reservoirs}(3) supplies at least $T$ internally disjoint candidate paths of length at most $H_0$ whose internal vertices lie in $C$.  Since fewer than $T$ vertices of $C$ have been used, and each used vertex lies internally on at most one path in this candidate family, at least one candidate avoids all previously used internal vertices.  Choose such a path.  Let $N_a\subseteq C$ be the union of the internal vertices of the paths chosen for $M_a$.

Finally set
\[
 B_a:=W_a\cup M_a\cup N_a
 \qquad(1\le a\le m).
\]
The sets $B_a$ are pairwise disjoint.  The subgraph $G[M_a\cup N_a]$ is connected, and every vertex of $W_a$ has a neighbor in $M_a$, so $G[B_a]$ is connected.  Since $W_1,\dots,W_m$ is a complete partition of $H'$ and $H'\subseteq G$, there is an edge of $G$ between $B_a$ and $B_b$ for every $a\ne b$.  Hence $B_1,\dots,B_m$ are branch sets of a $K_m$ minor of $G$.  The lower bound \eqref{eq:m-target} is exactly the asserted bound on $s$.
\end{proof}

\section{Proof of \cref{thomasonlower}}

Throughout this appendix, let $G$ be the random graph described in \cref{thomasonlower}, let $V = V(G)$, and let $(X_1, \ldots X_k)$ be a partition of $V$. For graphs \(F\) and \(H\), we write \(F\preccurlyeq H\) if \(F\) is a
minor of \(H\).  The \emph{Hadwiger number} of a nonempty graph \(H\) is
\[
  \had(H):=\max\{t\in\N:K_t\preccurlyeq H\}.
\]

The classical homogeneous permissible-partition argument appears in
Bollob\'as, Catlin and Erd\H{o}s~\cite{MR593989}.  Their calculation treats \(p=\tfrac12\), and the
paper also records the extension to every fixed \(0<p<1\)
\cite[p.~198]{MR593989}.  We give a complete proof of the inhomogeneous
estimate needed here.

\subsection{A finite probability estimate}

In
this section, no asymptotic assumptions are imposed.  Put
\[
  n_i:=|X_i|,\qquad
  \lambda_{ij}:=-\log(1-p_{ij}),
\]
and define the aggregate logarithmic edge weight
\begin{equation}\label{eq:def-Lambda}
  \Lambda
  :=\sum_{i,j\in[k]}\alpha_i\alpha_j\lambda_{ij}.
\end{equation}
Thus \(\lambda_{ij}=\lambda_{ji}\geq0\), and the quantity \(L\) in
\cref{thomasonlower} satisfies \(L=\sqrt{\Lambda}\).

An ordered partition
\[
  \mathcal A=(A_1,\dots,A_s)
\]
of \(V\) into nonempty sets is called \emph{permissible} for a graph
\(H\) if, for every \(1\leq a<b\leq s\), at least one edge of \(H\)
has one endpoint in \(A_a\) and the other in \(A_b\).  Connectivity of
the graphs \(H[A_a]\) is not required.\todo{maybe remove the connectivity sentence, but let's see if you find it helpful}

We first record a straightforward property of complete graph minors.

\begin{obs}
\label{obs:minor-partition}
Let \(H\) be a graph and let \(1\leq s\leq |V(H)|\).  If
\(K_s\preccurlyeq H\), then \(V(H)\) has a permissible ordered
partition into \(s\) nonempty parts.
\end{obs}

For \(v\in V\), let \(\tau(v)\) be the unique index for which
\(v\in X_{\tau(v)}\).  Give each potential edge \(\{u,v\}\) the
weight \(\lambda_{\tau(u),\tau(v)}\), and let
\[
  T:=\sum_{\{u,v\}\in\binom V2}
       \lambda_{\tau(u),\tau(v)}
\]
be the total weight of all potential edges.

\begin{lem}\label{lem:total-weight}
The following identity and inequality hold:
\begin{equation}\label{eq:total-weight}
  T
  =\frac{n^2\Lambda}{2}
   -\frac12\sum_{i=1}^k n_i\lambda_{ii}
  \leq\frac{n^2\Lambda}{2}.
\end{equation}
\end{lem}

\begin{proof}
By separating pairs whose vertices lie in the same class from those
whose vertices lie in different classes, and using
\(\lambda_{ij}=\lambda_{ji}\), we obtain
\begin{align*}
 T
 &=\sum_{i=1}^k\binom{n_i}{2}\lambda_{ii}
   +\sum_{1\leq i<j\leq k}n_in_j\lambda_{ij}\\
 &=\frac12\sum_{i,j\in[k]}n_in_j\lambda_{ij}
   -\frac12\sum_{i=1}^k n_i\lambda_{ii}.
\end{align*}
Since \(n_i=n\alpha_i\), the double sum in the last line is
\(n^2\Lambda\), proving the identity.  The inequality follows from
\(\lambda_{ii}\geq0\).
\end{proof}

\begin{lem}\label{lem:fixed-partition}
Let \(2\leq s\leq n\), and let
\(\mathcal A=(A_1,\dots,A_s)\) be a fixed ordered partition of \(V\)
into nonempty sets.  Then
\begin{equation}\label{eq:fixed-bound}
 \Prob(\mathcal A\text{ is permissible for }G)
 \leq
 \exp\left\{
  -\binom{s}{2}
   \exp\left(-\frac{n^2\Lambda}{s(s-1)}\right)
 \right\}.
\end{equation}
\end{lem}

\begin{proof}
For \(1\leq a<b\leq s\), set
\[
 W_{ab}:=
 \sum_{u\in A_a}\sum_{v\in A_b}
 \lambda_{\tau(u),\tau(v)}.
\]
By independence of the edge indicators and the definition of
\(\lambda_{ij}\),
\begin{align}
 \Prob\bigl(E_G(A_a,A_b)=\varnothing\bigr)
 &=\prod_{u\in A_a}\prod_{v\in A_b}
   (1-p_{\tau(u),\tau(v)})\notag\\
 &=\exp(-W_{ab}).
 \label{eq:no-edge}
\end{align}
For distinct unordered pairs \(\{a,b\}\), the events asserting the
existence of an edge between \(A_a\) and \(A_b\) are mutually independent. Equation \eqref{eq:no-edge} and the inequality \(1-x\leq e^{-x}\) give
\begin{align}
 \Prob(\mathcal A\text{ is permissible for }G)
 &=\prod_{1\leq a<b\leq s}
   \bigl(1-e^{-W_{ab}}\bigr)\notag\\
 &\leq
 \exp\left\{-\sum_{1\leq a<b\leq s}e^{-W_{ab}}\right\}.
 \label{eq:permissible-product}
\end{align}

Every potential edge whose endpoints belong to two different parts of
\(\mathcal A\) contributes to exactly one of the \(W_{ab}\), while
potential edges with both endpoints in one part contribute to none.
Lemma~\ref{lem:total-weight} therefore implies
\begin{equation}\label{eq:sum-W}
 \sum_{1\leq a<b\leq s}W_{ab}
 \leq T\leq\frac{n^2\Lambda}{2}.
\end{equation}
The function \(x\mapsto e^{-x}\) is convex. Writing \(M:=\binom{s}{2}\), Jensen's inequality and
\eqref{eq:sum-W} yield
\begin{align*}
 \sum_{1\leq a<b\leq s}e^{-W_{ab}}
 &\geq
 M\exp\left(-\frac1M\sum_{1\leq a<b\leq s}W_{ab}\right)\\
 &\geq
 M\exp\left(-\frac{n^2\Lambda}{2M}\right).
\end{align*}
Substituting this lower bound into
\eqref{eq:permissible-product} proves
\eqref{eq:fixed-bound}.
\end{proof}
\begin{cor}\label{cor:finite} For every integer \(s\) with \(2\leq s\leq n\),
\begin{equation}\label{eq:finite}
 \Prob\bigl(\had(G)\geq s\bigr)
 \leq
 s^n
 \exp\left\{
  -\binom{s}{2}
   \exp\left(-\frac{n^2\Lambda}{s(s-1)}\right)
 \right\}.
\end{equation}
\end{cor}

\begin{proof}
Every ordered partition of \(V\) into \(s\) nonempty parts is represented
by a surjection \(V\to[s]\).  Since there are \(s^n\) maps
\(V\to[s]\) altogether, there are at most \(s^n\) such ordered
partitions.  By Observation~\ref{obs:minor-partition}, the event
\(\{\had(G)\geq s\}\) is contained in the union, over these partitions,
of the events that the partition is permissible.  The union bound and
Lemma~\ref{lem:fixed-partition} now give \eqref{eq:finite}.
\end{proof}

\subsection{Uniform asymptotics}

We isolate the uniform estimate that converts
Corollary~\ref{cor:finite} into Theorem~\ref{thomasonlower}.

\begin{lem}\label{lem:uniform-decay}
Fix \(0<\eps<1\), and put
\[
 B:=-\log\eps.
\]
There exists \(n_1=n_1(\eps)\) such that the following holds.  If
\(n\geq n_1\), if
\begin{equation}\label{eq:Lambda-range}
 \eps^2<\Lambda\leq B,
\end{equation}
and if
\begin{equation}\label{eq:def-R-s}
 R:=(1+\eps)\frac{n\sqrt{\Lambda}}{\sqrt{\log n}},
 \qquad s_0:=\lceil R\rceil,
\end{equation}
then \(2\leq s_0\leq n\) and $G$ satisfies
\begin{equation}\label{eq:uniform-tail}
 \Prob\bigl(\had(G)\geq s_0\bigr)\leq\eps.
\end{equation}
\end{lem}

\begin{proof}
Set
\[
 \theta:=1+\frac{\eps}{2},
 \qquad
 \gamma:=1-\theta^{-2}>0,
 \qquad
 c:=1-\left(\frac{\theta}{1+\eps}\right)^2>0.
\]
The range \eqref{eq:Lambda-range} gives the uniform bounds
\begin{equation}\label{eq:R-bounds}
 (1+\eps)\eps\frac{n}{\sqrt{\log n}}
 <R
 \leq
 (1+\eps)\sqrt{B}\frac{n}{\sqrt{\log n}}.
\end{equation}
The left-hand side tends to infinity, while the right-hand side is
\(o(n)\).  We may therefore choose \(n_1\), initially, so large that
for every admissible \(\Lambda\),
\begin{equation}\label{eq:R-large-small}
 R\geq\max\{2,c^{-1}\}
 \qquad\text{and}\qquad
 R+1\leq n.
\end{equation}
It follows at once that \(2\leq s_0\leq n\).

Because \(s_0\geq R\) and \(s_0-1\geq R-1\), the definition of \(c\)
and \eqref{eq:R-large-small} imply
\begin{align}
 s_0(s_0-1)
 &\geq R(R-1)
   =R^2\left(1-\frac1R\right)\notag\\
 &\geq (1-c)R^2
   =\theta^2\frac{n^2\Lambda}{\log n}.
 \label{eq:rounding}
\end{align}
Consequently,
\begin{equation}\label{eq:two-consequences}
 \frac{n^2\Lambda}{s_0(s_0-1)}
 \leq\frac{\log n}{\theta^2},
 \qquad
 \binom{s_0}{2}
 \geq\frac{\theta^2n^2\Lambda}{2\log n}.
\end{equation}

Apply Corollary~\ref{cor:finite} with \(s=s_0\).  If
\(\Prob(\had(G)\geq s_0)=0\), then \eqref{eq:uniform-tail} is immediate.
Otherwise, taking logarithms in \eqref{eq:finite}, and then using
\(s_0\leq n\), \eqref{eq:two-consequences}, and
\(\Lambda>\eps^2\), gives
\begin{align}
 \log\Prob\bigl(\had(G)\geq s_0\bigr)
 &\leq n\log s_0
 -\binom{s_0}{2}
  \exp\left(-\frac{n^2\Lambda}{s_0(s_0-1)}\right)\notag\\
 &\leq n\log n
 -\frac{\theta^2n^2\Lambda}{2\log n}\,
  n^{-\theta^{-2}}\notag\\
 &\leq n\log n
 -\frac{\theta^2\eps^2}{2}
  \frac{n^{\,2-\theta^{-2}}}{\log n}\notag\\
 &=n\log n\left(
   1-\frac{\theta^2\eps^2}{2}
     \frac{n^\gamma}{(\log n)^2}
  \right).
 \label{eq:log-tail}
\end{align}
Since \(\gamma>0\),
\begin{equation}\label{eq:idk}
 \frac{n^\gamma}{(\log n)^2}\longrightarrow\infty.
\end{equation}
The final expression in \eqref{eq:log-tail} depends only on $n$ and $\eps$, so by \eqref{eq:idk}, it tends to
\(-\infty\) as $n$ tends to $\infty$.  Increasing \(n_1(\eps)\) if necessary, we can make
that expression at most \(\log\eps\) for every \(n\geq n_1\).
Exponentiating proves \eqref{eq:uniform-tail}.
\end{proof}

\begin{proof}[Proof of Theorem~\ref{thomasonlower}]
If \(\eps\geq1\), the hypotheses are inconsistent: when \(\eps>1\)
there is no nonnegative \(p_{ij}\leq1-\eps\), while for \(\eps=1\)
all \(p_{ij}\) would be zero and hence \(L=0\).  Thus the assertion is
vacuous.  Assume henceforth that \(0<\eps<1\).

For every \(i,j\), the bound \(0\leq p_{ij}\leq1-\eps\) gives
\[
 0\leq\lambda_{ij}=-\log(1-p_{ij})\leq-\log\eps=B.
\]
Since
\[
 \sum_{i,j\in[k]}\alpha_i\alpha_j
 =\left(\sum_{i=1}^k\alpha_i\right)^2=1,
\]
equation \eqref{eq:def-Lambda} implies
\[
 0\leq\Lambda\leq B.
\]
Moreover, \(L=\sqrt{\Lambda}>\eps\), so
\(\Lambda>\eps^2\).  We may therefore apply
Lemma~\ref{lem:uniform-decay}.  For every \(n\geq n_1(\eps)\), it gives
\[
 \Prob\left(
  \had(G)\geq
  \left\lceil(1+\eps)\frac{nL}{\sqrt{\log n}}\right\rceil
 \right)
 \leq\eps.
\]

The Hadwiger number is an integer.  Hence the complementary event is
exactly the event that \(G\) has no \(K_s\) minor for any integer
\[
 s\geq(1+\eps)\frac{nL}{\sqrt{\log n}}.
\]
This event has probability at least \(1-\eps\), as required.  Taking \(n_0(\eps):=n_1(\eps)\) proves the theorem.
\end{proof}

\begin{remark}\label{rem:homogeneous}
If \(p_{ij}=p\) for all \(i,j\), then
\[
 \Lambda=-\log(1-p),\qquad
 L=\sqrt{-\log(1-p)}.
\]
Writing \(b=(1-p)^{-1}\), the threshold becomes
\[
 (1+\eps)\frac{n}{\sqrt{\log_b n}},
\]
which is the classical Bollob\'as--Catlin--Erd\H{o}s upper-bound scale.
\end{remark}

\newpage

\textbf{Sketch of 4-colour example} \\

To verify: this is simply by messing around with the statement of Thomason's theorem (Theorem 1 in this document). There could still be other details to work out.

Randomly four-colouring edges of a complete graph gives a largest minor of order $\sqrt{-\log\s{3/4}}n_* \approx 0.5364 n_*$. We discuss a four-colouring of a complete graph that gives a largest minor of order approximately $0.4704 n_*$.

We split a graph on $n$ vertices into 4 equal (or almost equal) components, with $n$ sufficiently large. Inside each component, edges take each of colours 1, 2, 3 with probability $1/3$. Any edge between components takes colour 4 with probability $1-p$, with $0 < p < 1$ to be chosen later. An edge between components 1 and 2 or between components 3 and 4 takes colour 1 with probability $p$. Similarly, an edge between components 1 and 3 or between components 2 and 4 takes colour 2 with probability $p$. Finally, an edge between components 1 and 4 or between components 2 and 3 takes colour 3 with probability $p$.  Essentially, the colour chosen with probability $p$ between each pair of components imitates 3 disjoint perfect matchings on 4 vertices, with only colour 4 going between all components.

From Theorem 1, the largest colour 4 minor has order $\sqrt{-\frac{12}{16}\ln(p)} n_*$. The remaining colours are symmetric, so we analyze colour 1 WLOG. The largest colour 1 minor can only span across 2 components. Applying Theorem 1 to only components 1 and 2, the largest colour 1 minor has order $\sqrt{-\frac2{16} \ln(2/3)-\frac2{16} \ln(1-p)} n_*$ Minimizing the maximum of these two quantities, for $0 < p < 1$, we get $$\max\{\sqrt{-\frac34\ln(p)}, \sqrt{-\frac18 \ln(2/3) - \frac18 \ln(1-p)}\} \approx 0.4704$$ when $p \approx 0.7445$.

Note that this construction does not yield a better result if we allow colour 4 edges inside a component.

\end{document}